%% ============================================================================
%% Chapter 17: Deployment & Infrastructure: Containers, Gateways & Edge
%% Mastering APIs: From Zero to Production Guru
%% ============================================================================
\chapter{Deployment \& Infrastructure: Containers, Gateways \& Edge}
\label{ch:deployment}

%% ----------------------------------------------------------------------------
%% Executive Summary
%% ----------------------------------------------------------------------------
Every chapter before this one ended at the same invisible boundary: the
API worked \emph{on your machine}. This chapter is about the far harsher
question of whether it keeps working while the machine underneath it is
being replaced, scaled, drained, upgraded, or killed. Deployment is not a
soup of YAML that ops people sprinkle on finished code; it is the
continuation of the API contract by other means. A client does not
experience your code --- it experiences the behavior of your code
\emph{while infrastructure is changing around it}. The difference between
a platform and a pile of services is whether a routine Tuesday-afternoon
deploy is indistinguishable, from the caller's side, from no deploy at all.

Three ideas organize everything that follows. First, the \emph{immutable
artifact}: a container image is the unit of deployment, and its quality is
a security and reliability property, not a cosmetic one. Multi-stage
builds, pinned digests, non-root users, and read-only filesystems are not
hygiene theater; they shrink the attack surface to a fraction of a
convenience image and make every byte in production traceable to a
commit. Second, \emph{process discipline}: inside the container your
process is PID~1, which makes signal handling your responsibility, and
termination is a four-act protocol --- stop accepting, drain in-flight
work, time out, exit --- that must be choreographed with the platform's
\emph{preStop} hooks, \texttt{terminationGracePeriodSeconds}, and the
embarrassingly long lag between a Pod dying and the rest of the cluster
noticing. Third, \emph{progressive delivery}: you never actually deploy a
version; you migrate traffic between versions, gated by metrics, with a
rollback path rehearsed before it is needed~\cite{nygard2018releaseit,beyer2016sre}.

We treat the API gateway as what it operationally is --- a policy
enforcement point that absorbs cross-cutting concerns (JWT validation,
rate limiting, request transformation, canary routing) so the service
code does not --- and we draw the boundary lines between gateway, ingress,
and service mesh precisely, because muddling them produces double TLS
terminations, duplicated rate limits, and unowned failures. On Kubernetes
we go past the tour and into the semantics that bite: the exact failure
behavior of startup versus readiness versus liveness probes (demonstrated
with a deterministic simulator, not asserted), the CPU-throttling latency
trap that \texttt{limits} create, HPA stabilization windows,
PodDisruptionBudgets, and the rolling-update arithmetic of
\texttt{maxSurge} and \texttt{maxUnavailable}. We close the loop at the
edge: TLS termination placement, certificate automation, CDN caching rules
for APIs, WAF essentials, and the north-south/east-west distinction that
decides which of these boxes a packet even visits.

Everything runnable is offline, deterministic, and localhost-only: a
graceful-shutdown proof that catches a termination signal \emph{during} a
slow request and shows the drain completing; a probe-semantics simulator
that replays a boot, a dependency outage, and a deadlock against the
kubelet's rules; and a canary analyzer that runs Wald's sequential
probability ratio test over synthetic metric series and promotes or rolls
back on evidence, not vibes. The Dockerfile and the Kubernetes manifests
are complete and production-shaped but are \emph{reviewed line by line},
not applied --- reading infrastructure critically is the skill this
chapter exists to teach.

%% ----------------------------------------------------------------------------
\section{Learning Objectives \& Prerequisites}
\label{sec:ch17-objectives}

\subsection*{Learning Objectives}
After completing this chapter, its lab, and its exercises, you will be able
to:

\begin{enumerate}[label=\textbf{LO-\arabic*.}, leftmargin=2.6cm]
  \item Author a production multi-stage Dockerfile for a FastAPI service:
        separate build and runtime stages, choose between slim and
        distroless bases, run as a non-root \texttt{USER}, control the
        build context with \texttt{.dockerignore}, order instructions for
        layer-cache reuse, pin base images by digest, and wire
        \texttt{HEALTHCHECK} and read-only filesystems --- and justify
        every line as attack-surface or reproducibility engineering.
  \item Explain the PID~1 problem and implement the graceful-shutdown
        sequence --- stop accepting, drain, timeout, exit --- including
        uvicorn's \texttt{SIGTERM} behavior, \emph{preStop} hooks, and
        load-balancer connection draining, and prove drain semantics
        empirically with the chapter's signal-driven demo.
  \item Compare Kong, Envoy, and NGINX as API gateways across routing,
        authentication offload, rate limiting, transformation, and canary
        support; position gateway versus ingress versus service mesh; and
        argue for the gateway as the policy enforcement point for
        north-south traffic.
  \item Design Kubernetes workloads for APIs with correct probe semantics
        (startup gates liveness/readiness; readiness removes endpoints
        without restarts; liveness restarts only process-local death),
        resource requests and limits that avoid throttle-induced latency,
        HPA metrics and stabilization windows, PodDisruptionBudgets, and
        rolling-update parameters chosen from replica-count arithmetic.
  \item Select and operate a progressive-delivery strategy --- blue-green,
        canary with traffic-splitting math and metric-gated promotion, or
        feature flags --- including automated rollback signals and the
        expand-and-contract database discipline that makes rollbacks
        possible at all (Chapter~\ref{ch:versioning}).
  \item Place TLS termination deliberately (edge, re-encrypt,
        passthrough), automate certificate lifecycles, define CDN cache
        behavior for APIs (what to cache, cache keys, purge semantics),
        apply WAF essentials, and classify any hop as north-south or
        east-west.
\end{enumerate}

\subsection*{Prerequisites}
This chapter assumes Chapters~0--16, with four load-bearing dependencies.
Chapter~\ref{ch:fastapi} (Building APIs with FastAPI) supplies the Orders
API we containerize and the \texttt{TestClient}/uvicorn mechanics the
shutdown demo leans on~\cite{fastapidocs}. Chapter~\ref{ch:rate_limiting}
(Rate Limiting) supplies the algorithms the gateway offloads and the
header contract the edge must preserve.
Chapter~\ref{ch:microservices} (Microservices \& Resilience) supplies
timeouts, retries, and circuit breakers --- the behaviors that make
connection draining and endpoint-removal lag survivable rather than
fatal~\cite{newman2021microservices}. Chapter~\ref{ch:api_security} (API
Security Engineering) supplies the JWT validation, secret handling, and
transport-security vocabulary reused at the gateway and the edge. The
expand-and-contract migration discussion calls back to
Chapter~\ref{ch:versioning} (Versioning \& Evolution). The Twelve-Factor
App's process and config factors~\cite{twelvefactor} are referenced
throughout; a passing familiarity with the release-engineering chapter of
the Google SRE book~\cite{beyer2016sre} helps but is not required.

\subsection*{Setup}
All runnable listings execute offline on Windows with Python~3.11+ and
localhost only. One-time setup from PowerShell in a scratch directory:

\begin{minted}{text}
python -m venv .venv
.\.venv\Scripts\Activate.ps1
pip install fastapi uvicorn httpx pytest pyyaml
\end{minted}

The Dockerfile and Kubernetes manifests are reviewed, not applied, so no
Docker daemon or cluster is required. If you want to go further,
Section~\ref{sec:ch17-deps} gives \texttt{winget} install lines for the
optional toolchain (\texttt{docker}, \texttt{kubectl}, \texttt{k6}).

%% ----------------------------------------------------------------------------
\section{Deep Technical Mechanics \& Architectural Concepts}
\label{sec:ch17-mechanics}

\subsection{The Immutable Artifact: What a Container Image Actually Is}
\label{sec:ch17-image}

A container image is not a lightweight virtual machine and it is not a
zip file of your app. It is a \emph{content-addressed stack of filesystem
layers} plus a small JSON manifest: each layer is a diff, each diff is
addressed by the SHA-256 of its contents, and the image itself is
addressed by the digest of the manifest. Two operational facts follow
immediately, and both are security facts before they are convenience
facts. First, layers are shared and cached: two images built from the
same base literally share those bytes on disk, which is why base-image
choice dominates image size and why instruction ordering dominates build
time. Second, a \emph{tag} like \texttt{3.12-slim} is a mutable pointer a
registry can retarget at any time, while a \emph{digest}
(\texttt{python@sha256:\ldots}) names exactly one manifest forever.
Pinning by digest is the difference between ``we deployed what we
tested'' and ``we deployed whatever the tag pointed at that morning.''

\begin{definition}[Immutable artifact]
An \emph{immutable artifact} is a deployable unit whose contents are fixed
at build time and identified by a cryptographic digest, so that every
running instance can be traced to exactly one build, and a rollback is a
pointer change to a previous digest rather than a new build.
\end{definition}

\textbf{Multi-stage anatomy.} A multi-stage Dockerfile contains at least
two \texttt{FROM} instructions. The first stage --- the \emph{build}
stage --- is allowed to be fat: compilers, headers, package managers, and
your full dependency resolver live there, because that stage is discarded.
The final stage --- the \emph{runtime} stage --- receives only copied
artifacts: a virtual environment or a \texttt{site-packages} tree, your
source, nothing else. The attack-surface math is simple and brutal: every
binary in the runtime image is a thing an attacker with remote code
execution can invoke, so a runtime stage containing \texttt{gcc},
\texttt{curl}, and a shell is a gift; a runtime stage containing
\texttt{python}, your wheels, and no shell at all forces an attacker to
bring their own tooling through whatever hole they already have.

\textbf{Base-image spectrum.} The practical spectrum runs from
\texttt{python:3.12} (full Debian, $\sim$1~GB, everything present,
everything exposable) through \texttt{python:3.12-slim}
($\sim$150~MB, no compiler, a shell and coreutils remain) to
\textit{distroless}
(\texttt{gcr.io/distroless/python3}: Python, glibc, CA certificates, and
\emph{nothing else} --- no shell, no package manager, typically
$\sim$50--80~MB for a FastAPI service). Distroless is the strongest
default for interpreters that need no OS packages; slim is the pragmatic
choice when you need one or two system libraries or when your
organization's vulnerability scanner needs an OS package database to
reason about. The wrong default is the full image ``because something
might need it'' --- something might also exploit it.

\textbf{Non-root, read-only, health-checked.} Three more lines of the
same argument. Containers on Docker and Kubernetes run as UID~0 by
default; a process that escapes as container-root is far closer to host
root than one that first has to escalate from an unprivileged UID, so the
runtime stage ends with \texttt{USER} set to a created non-root account.
A \texttt{readOnlyRootFilesystem} (or \texttt{docker run --read-only})
turns whole exploit classes --- webshell drops, config tampering, binary
replacement --- into no-ops, at the price of explicitly mounting writable
\texttt{tmpfs} for the few paths that genuinely need writes.
\texttt{HEALTHCHECK} (or its Kubernetes successor, the probe, which
subsumes it there) encodes the service's own opinion of its health into
the artifact instead of leaving it tribal.

\textbf{Layer-caching discipline.} BuildKit caches each instruction's
result and invalidates the cache from the first changed instruction
downward. The discipline is therefore: order instructions from
least-frequently-changed to most-frequently-changed (base image, system
packages, dependency manifests, dependency install, \emph{then} source),
and control what \texttt{COPY . .} can even see with
\texttt{.dockerignore} --- excluding \texttt{.git}, virtual environments,
test artifacts, and any secrets that wandered into the working tree. An
unignored \texttt{.env} copied into an image layer is a credential leak
that survives deletion, because the layer remains in the image history.

\begin{warningbox}
A file deleted in a later \texttt{RUN rm} is still present in the earlier
layer and recoverable from the image. ``Install, use, delete in one
\texttt{RUN}'' is a size optimization; ``copy secret, delete secret in a
later layer'' is a breach. Secrets never enter the build context or a
layer --- use BuildKit \texttt{--mount=type=secret} for build-time needs
and runtime secret mounts for deploy-time needs
(Chapter~\ref{ch:api_security}).
\end{warningbox}

\subsection{Process Discipline: PID 1, Signals, and the Four-Act Shutdown}
\label{sec:ch17-process}

In a Linux namespace, PID~1 is special in two ways that matter to an API
author. First, orphaned processes are re-parented to PID~1, which is
responsible for reaping zombies; second, and more important here, a
process only receives a signal it has not installed a handler for if that
signal's default disposition applies --- and for PID~1 the kernel
\emph{only delivers signals for which a handler is explicitly installed}.
An application that spawns a shell wrapper that spawns the server, or a
naive process that never calls \texttt{signal.signal}, can therefore be
unreachable by \texttt{SIGTERM}: Kubernetes asks politely, waits
\texttt{terminationGracePeriodSeconds}, and then issues
\texttt{SIGKILL} --- every in-flight request severed mid-write. The fixes
are boring and absolute: use \texttt{exec} (or no shell at all --- the
exec form of \texttt{ENTRYPOINT}/\texttt{CMD}) so your server \emph{is}
PID~1, and let the server framework install proper handlers. uvicorn
does: on \texttt{SIGTERM} it sets \texttt{should\_exit}, stops accepting
new connections, waits for in-flight requests up to
\texttt{timeout\_graceful\_shutdown}, runs ASGI lifespan shutdown, and
exits.

\begin{definition}[Graceful shutdown sequence]
\emph{Graceful shutdown} is the ordered protocol: (1)~\textbf{stop
accepting} --- close listeners and fail readiness so no new traffic is
routed; (2)~\textbf{drain} --- allow in-flight requests to finish up to a
budget; (3)~\textbf{timeout} --- abort stragglers that exceed the budget
so the process cannot hang forever; (4)~\textbf{exit} --- release
resources and terminate with a status the supervisor can interpret.
Skipping act~(1) routes new work to a dying instance; skipping act~(2)
drops requests clients paid for; skipping act~(3) turns a deploy into a
wedged rollout.
\end{definition}

The platform choreography has two sharp edges. First, the
\textbf{preStop hook}: when Kubernetes decides to terminate a Pod, it
runs the container's \texttt{preStop} hook \emph{before} sending
\texttt{SIGTERM}, and endpoint removal (updating the Service's Endpoint
objects, propagating to kube-proxy/iptables or the gateway's load
balancer) happens \emph{concurrently} --- not before, not after. A Pod
can therefore receive \texttt{SIGTERM} while stale endpoints still route
it fresh traffic for several seconds. The standard mitigation is a
deliberate \texttt{sleep} in \texttt{preStop} (5--10~seconds is common)
so the process keeps serving while the endpoint update fans out, then
accepts the signal into a quieted listener. Second,
\textbf{connection draining at the load balancer} is the same idea one
tier up: cloud LBs and gateways mark the target ``draining,'' stop
sending new connections, and wait for existing ones --- your
\texttt{terminationGracePeriodSeconds} must exceed the LB drain timeout
plus your slowest request, or the platform kills the process with work
still in flight. Section~\ref{sec:ch17-lab} replays the war story in
which exactly one missing \texttt{preStop} sleep turned a routine rolling
deploy into dropped checkouts.

\begin{keyconceptbox}
Termination is a \emph{race you must rig}, not an event you can order:
SIGTERM, endpoint removal, and LB draining start concurrently. You rig it
by making your server tolerant of late traffic (preStop sleep), tolerant
of early signals (drain logic), and explicit about the budget
(terminationGracePeriodSeconds $>$ slowest request $+$ preStop delay).
\end{keyconceptbox}

\subsection{The API Gateway: Policy Enforcement Point for North-South Traffic}
\label{sec:ch17-gateway}

An API gateway is a reverse proxy with opinions. It terminates client
connections at the edge of your platform and applies the cross-cutting
concerns that must be uniform across services: authentication offload,
rate limiting, request/response transformation, routing, and telemetry.
The architectural argument is the one from Chapter~\ref{ch:microservices}:
any policy implemented in $N$ services is implemented $N$ times,
differently, and audited never; a policy enforced once at the gateway is
uniform, versioned, and testable. \textbf{JWT validation at the edge} is
the canonical example --- the gateway verifies the signature, expiry, and
issuer against cached JWKS and rejects forged or expired tokens before
they consume a single cycle of service CPU, while services receive an
already-authenticated identity in a trusted header
(Chapter~\ref{ch:api_security}). \textbf{Rate limiting at the gateway}
(Chapter~\ref{ch:rate_limiting}) protects every upstream uniformly with
one configuration language instead of $N$ middleware copies.

\begin{table}[h]
  \centering
  \small
  \begin{tabularx}{\textwidth}{l X X X}
    \toprule
    \textbf{Capability} & \textbf{Kong} & \textbf{Envoy} & \textbf{NGINX (OSS / Plus)} \\
    \midrule
    Core proxy engine & NGINX + OpenResty (Lua) & purpose-built C++ L4/L7 proxy & NGINX \\
    Routing & rich: host/path/method/headers, plugins & very rich: L7 match, weights, retries & host/path/regex map; headers via config \\
    Auth offload (JWT) & first-class plugin (JWKS, claims checks) & JWT authn filter (JWKS, claims) & Plus native; OSS via njs/Lua modules \\
    Rate limiting & plugin (local/Redis, per consumer/route) & local token bucket + global rate-limit service & \texttt{limit\_req} (GCRA); Plus adds key-value store \\
    Request transformation & plugin ecosystem (Lua) & Lua/Wasm filters; header/body mutation & \texttt{sub\_filter}, headers, njs scripting \\
    Canary routing & traffic-split plugin, weighted upstreams & native weighted clusters (first-class) & \texttt{split\_clients} / Plus API \\
    Extensibility model & plugin hub + custom Lua/Go/JS/Python & Wasm/Lua/ext\_proc; the mesh data-plane standard & modules at compile time; njs runtime \\
    Config style & declarative (DB or deck YAML), dynamic API & xDS APIs (built for control planes) & static files + reload (Plus: dynamic API) \\
    Operational posture & gateway-first, admin API, enterprise suite & building block: gateways/meshes embed it & ubiquitous, minimal deps, least dynamic \\
    \bottomrule
  \end{tabularx}
  \caption{Gateway capabilities matrix. Envoy is the substrate many
  gateways and meshes (Istio, Consul, Contour, Gloo) are built from; Kong
  is a batteries-included gateway on NGINX; NGINX itself is the minimal,
  static-configuration baseline.}
  \label{tab:ch17-gateway-matrix}
\end{table}

\textbf{Gateway vs.\ ingress vs.\ mesh.} These three are conflated
constantly and are distinguished by \emph{what they terminate and where
their policy boundary sits}. \emph{Ingress} (the Kubernetes resource and
its controllers) is the narrowest: a cluster-standard way to route
north-south HTTP(S) to Services, with TLS termination and host/path
rules --- a gateway subset with a portable vocabulary. An \emph{API
gateway} is a product-tier edge: everything ingress does plus auth
offload, rate limiting, transformation, developer portals, and
analytics --- the \textbf{policy enforcement point} for north-south
traffic. A \emph{service mesh} (sidecar or ambient proxies on every
workload) governs \emph{east-west} traffic: service-to-service mTLS,
retries, circuit breaking, and telemetry \emph{inside} the platform. The
common production shape is all three in sequence: CDN/WAF at the
provider edge, gateway at the platform edge, mesh between services ---
and the failure mode to avoid is enforcing the same concern in two tiers
(e.g., rate limiting at gateway \emph{and} mesh with different budgets)
so neither is the authoritative budget.

\subsection{Kubernetes for APIs: The Semantics That Bite}
\label{sec:ch17-k8s}

\textbf{Anatomy in one paragraph.} A \emph{Deployment} declares desired
state --- $N$ identical Pods from one Pod template --- and a controller
continuously reconciles reality toward it. A \emph{Pod} is the scheduling
unit: one or more containers sharing a network namespace (one IP) and
volumes. A \emph{Service} is a stable virtual IP and DNS name that load
balances across the Pods matching its selector, via an endpoints
controller that tracks \emph{ready} Pods. Everything below is about the
gaps between these declarations --- the gaps are where requests die.

\textbf{Probes: three questions, three consequences.} The kubelet asks
your container three distinct questions, and the failure consequence of
each is different --- this is the most cargo-culted area of Kubernetes,
so we state it as a table and then \emph{prove} it with the simulator in
Section~\ref{sec:ch17-code}.

\begin{table}[h]
  \centering
  \small
  \begin{tabularx}{\textwidth}{l X X X}
    \toprule
    & \textbf{Startup} & \textbf{Readiness} & \textbf{Liveness} \\
    \midrule
    Question & Has the app finished initializing? & Should this Pod receive traffic \emph{right now}? & Is the process unrecoverably stuck? \\
    On failure & Retried within budget (\texttt{period} $\times$ \texttt{failureThreshold}); budget exhausted $\Rightarrow$ container \textbf{restarted} & Pod \textbf{removed from Service endpoints}; no restart, no reschedule; re-added on first success & Failure streak $\Rightarrow$ container \textbf{restarted in place} (same Pod, new container) \\
    Gates & Disables liveness \& readiness until first success & Nothing; evaluated only after startup succeeds & Nothing; evaluated only after startup succeeds \\
    Should fail when & Still booting & A \emph{recoverable}/external problem: dependency down, overloaded, draining & A \emph{process-local} death: deadlock, wedged event loop \\
    Typical check & \texttt{/healthz/startup} (cheap, no deps) & \texttt{/healthz/ready} (may check critical deps, shallowly) & \texttt{/healthz/live} (no deps --- deadlocks must be the only failure) \\
    \bottomrule
  \end{tabularx}
  \caption{Probe semantics. The load-bearing row is the second:
  readiness failure changes \emph{routing}, liveness failure changes
  \emph{process state}. A readiness probe that also gates liveness
  restarts Pods for other services' outages; a liveness probe that
  checks dependencies turns a database blip into a fleet-wide restart
  storm.}
  \label{tab:ch17-probes}
\end{table}

The \textbf{startup probe} exists because boot time and liveness budget
are unrelated quantities. Without it, a slow-starting application (JVM
warmup, model load, cache prime) is indistinguishable from a deadlocked
one, and the liveness budget becomes a boot deadline you never meant to
set; the simulator's Scenario~B crashes exactly this way. The
\textbf{readiness/liveness split} encodes a recovery theory: readiness
says ``wait'' (the problem may be elsewhere --- restarting me cannot fix
the database), liveness says ``kill me'' (the problem \emph{is} me).
\textbf{Endpoint-removal lag} completes the picture: readiness and
termination only update the endpoints object; kube-proxy programming,
DNS caches, and gateway target lists converge seconds later, which is
why Section~\ref{sec:ch17-process} insists on the preStop sleep and why
clients still need retries (Chapter~\ref{ch:microservices}).

\begin{figure}[h]
\centering
\begin{tikzpicture}[font=\small]
  % timeline axis
  \draw[->, thick] (0,0) -- (13.2,0) node[right, font=\footnotesize] {time};
  % phase bands
  \fill[packtorange!12] (0,0) rectangle (3.4,2.6);
  \fill[blue!6] (3.4,0) rectangle (13.0,2.6);
  \node[note] at (1.7,2.9) {startup phase};
  \node[note] at (8.2,2.9) {steady state (startup probe retired)};
  % probe rows
  \node[flowstep, anchor=west] at (0.2,2.0) {startup probe};
  \node[flowstep, anchor=west] at (0.2,1.1) {readiness probe};
  \node[flowstep, anchor=west] at (0.2,0.35) {liveness probe};
  % startup probe events
  \foreach \x in {0.9,1.7,2.5} \draw[red!70,thick] (\x,1.85) -- (\x,2.15);
  \draw[green!50!black,very thick] (3.3,1.85) -- (3.3,2.15) node[above, note, align=center, yshift=2pt] {first success};
  % readiness events
  \foreach \x in {4.2,5.0,7.4} \draw[green!50!black,thick] (\x,0.95) -- (\x,1.25);
  \foreach \x in {5.8,6.6} \draw[red!70,thick] (\x,0.95) -- (\x,1.25);
  \draw[green!50!black,very thick] (8.2,0.95) -- (8.2,1.25);
  \node[note, align=left] at (9.6,1.1) {3 fails $\Rightarrow$ endpoint \emph{removed};\\1 success $\Rightarrow$ re-added};
  % liveness events
  \foreach \x in {4.2,6.6,8.2} \draw[green!50!black,thick] (\x,0.2) -- (\x,0.5);
  \foreach \x in {9.8,10.8,11.8} \draw[red!70,thick] (\x,0.2) -- (\x,0.5);
  \draw[->, thick, red!70] (12.1,0.35) -- (12.8,0.35) node[right, note, align=left] {3 fails $\Rightarrow$\\restart};
  % annotations
  \node[note, align=center] at (1.8,1.55) {fails while booting\\(budget: period$\times$threshold)};
  \draw[dashed, packtgray] (3.4,0) -- (3.4,2.6);
\end{tikzpicture}
\caption{Probe lifecycle. The startup probe alone is evaluated during
boot, protecting slow starts from the liveness budget. After its first
success, readiness (routing consequence) and liveness (restart
consequence) run on their own periods.}
\label{fig:ch17-probes}
\end{figure}

\textbf{Resources: requests schedule, limits throttle.} Kubernetes uses
\texttt{requests} for scheduling (bin-packing: a node must have the
requested CPU/memory free) and as the QoS input, and \texttt{limits} as
hard ceilings with two very different enforcement mechanisms. Exceed the
memory limit and the kernel OOM-kills the container --- violent, visible,
logged. Exceed the CPU \emph{limit} and nothing fails: the kernel's
Completely Fair Scheduler \emph{throttles} the cgroup for the remainder
of each 100~ms quota period. The pitfall is that throttling is quantized:
a service with 8 worker threads and a 400m CPU limit can consume its
40~ms-of-quota-per-period in the first 5~ms of a burst, then freeze every
thread --- including the one holding a request --- for up to 95~ms. The
symptom is p99 latency inflated by \emph{exactly} multiples of ~100~ms
while average CPU sits far below the limit, which is why CPU-limit
throttling is the classic cause of ``mystery tail latency'' in Go and
Python services alike. Operational guidance: always set CPU
\emph{requests} (they also weight CFS shares under contention), set
memory requests~=~limits (memory cannot be throttled, so a limit without
a matching request is pure risk), and either omit CPU limits on
latency-sensitive services or set them generously above p99 usage and
alert on \texttt{container\_cpu\_cfs\_throttled\_periods\_total}.

\textbf{HPA: metrics, windows, and flapping.} The Horizontal Pod
Autoscaler computes $\text{desired} = \lceil \text{current} \times
\text{currentMetric} / \text{targetMetric} \rceil$ every sync period from
resource metrics (CPU/memory utilization against \emph{requests}),
custom metrics (requests per second per pod), or external metrics
(queue depth). Two details separate operators from tourists: utilization
is computed against \emph{requests}, so a service with inflated requests
never autoscales no matter how hot it runs; and the default
\texttt{scaleDown} stabilization window (300~s) exists because scaling
down kills in-flight connections --- the window makes downscale
decisions from the \emph{highest} recommendation in the window, trading
money for connection stability. Downscale \emph{rate} limits
(\texttt{policies}: e.g., remove at most 1 pod or 10\% per 60~s) are the
second half of the same idea.

\begin{table}[h]
  \centering
  \small
  \begin{tabularx}{\textwidth}{l l l X}
    \toprule
    \textbf{\texttt{maxSurge}} & \textbf{\texttt{maxUnavailable}} & \textbf{Serving floor ($R{=}4$)} & \textbf{Cost / risk} \\
    \midrule
    1 & 0 & 4 & one surge pod's quota; zero capacity loss --- the default for user-facing APIs \\
    0 & 1 & 3 & no extra quota; 25\% capacity dip per rollout wave \\
    1 & 1 & 3 & fastest rollout; both costs at once \\
    0 & 0 & --- & cannot roll at all; a copy-paste wedge, not a strategy \\
    \bottomrule
  \end{tabularx}
  \caption{Rolling-update arithmetic for a 4-replica Deployment.
  Percentages are allowed for both fields (rounded: surge up,
  unavailable down), which keeps the policy correct as HPA moves $R$.}
  \label{tab:ch17-rollout}
\end{table}

\textbf{PDBs and rolling updates.} A PodDisruptionBudget states how many
Pods of a Service may be \emph{voluntarily} unavailable
(\texttt{minAvailable} or \texttt{maxUnavailable}); node drains and
cluster upgrades respect it, crashes do not (they are involuntary).
Without a PDB, \texttt{kubectl drain} will happily evict every replica
of your API simultaneously. Rolling updates trade capacity for safety via
two parameters: \texttt{maxSurge} (extra Pods above desired allowed
during the rollout --- needs quota headroom) and
\texttt{maxUnavailable} (desired Pods allowed to be down). With $R{=}4$
replicas, \texttt{maxSurge: 1, maxUnavailable: 0} keeps serving capacity
at $\ge 4$ throughout at the cost of one extra pod's resources; with
\texttt{maxUnavailable: 1} you never exceed $R$ but serve a window at
$R{-}1$. The interaction that bites: HPA plus a bad PDB
(\texttt{minAvailable} equal to replicas) deadlocks drains, and a
readiness gate that never passes (new version's probe path renamed)
turns \texttt{maxUnavailable: 0} into a rollout that cannot proceed ---
which is the safe failure, and why \texttt{progressDeadlineSeconds}
exists to fail the rollout loudly instead of wedging it silently.

\subsection{Progressive Delivery: Migrating Traffic, Not Versions}
\label{sec:ch17-progressive}

\textbf{Blue-green} keeps two full environments; the cutover is an atomic
traffic switch (DNS, LB target group, or Service selector), and rollback
is the same switch flipped back. Its virtues are instant rollback and a
fully warm target; its costs are 2x infrastructure and the cold-start
problem merely moved (the idle environment rots unless exercised).
\textbf{Canary} releases the new version to a fraction of traffic and
promotes through steps --- 1\%, 5\%, 25\%, 50\%, 100\% is a common
ladder --- comparing SLIs at each step. The traffic-splitting math has
two subtleties. First, error budgets are ratios, and ratios need volume:
at a 0.4\% baseline error rate, a 1\% canary serving 100~requests shows
zero errors with probability $\approx 67\%$ even if the canary is
perfectly healthy --- small slices are statistically blind, so canary
steps must be sized by \emph{requests}, not just percentages, which is
why the analyzer in Section~\ref{sec:ch17-code} works in fixed-size
batches. Second, weighted splitting must be \emph{stable} for
stateful-feeling flows: session-affinity-aware hashing (on API key or
tenant) prevents one client oscillating between versions mid-workflow,
which matters whenever response shapes differ
(Chapter~\ref{ch:versioning}).

\begin{definition}[Metric-gated promotion]
A \emph{metric-gated promotion} advances traffic to the next canary step
only when a pre-registered analysis over a fixed window passes --- and
rolls back automatically when regression evidence crosses a threshold.
The gate is a statistical decision procedure (here: Wald's SPRT with
error budgets $\alpha$ for false-promote and $\beta$ for false-rollback),
not a human watching a dashboard; the human designed the procedure.
\end{definition}

\textbf{Automated rollback signals} worth wiring: error-rate delta
versus baseline (the SPRT in Listing~\ref{lst:ch17-canary}), latency
quantile guardrails (p95/p99 against a hard factor of baseline), and
saturation signals (CPU throttling, connection-pool exhaustion). Signals
that \emph{look} useful and are not: average latency (hides tails),
absolute error counts (denominator-blind), and business metrics on short
windows (too noisy to gate on minute timescales).

\textbf{Feature flags vs.\ deployment strategies} are complements, not
rivals: deployments change \emph{which binary} runs, flags change
\emph{which code paths} execute inside one binary. Flags decouple release
from deploy (ship dark, enable per-tenant, kill-switch instantly) at the
price of combinatorial code paths and flag-debt; canaries decouple blast
radius from deploy speed at the price of two contract versions in flight
simultaneously. Both, therefore, rest on the same foundation:
\textbf{database compatibility}. The expand-and-contract discipline of
Chapter~\ref{ch:versioning} is what makes any of this safe: migrations
run in the \emph{expand} step (add column/table, never rename or drop),
both binary versions tolerate both shapes, and \emph{contract} (drop,
rename) ships only after every old binary is gone. Deploying a
destructive migration \emph{before} the compatible code is the one
rollout error rollback cannot fix: the old binary comes back, the column
it needs does not. We name this rule \emph{migrate-expand, code,
migrate-contract} and treat violations as release blockers.

\begin{figure}[h]
\centering
\begin{tikzpicture}[font=\small]
  % axes
  \draw[->, thick] (0,0) -- (13.4,0) node[right, font=\footnotesize] {time};
  \draw[->, thick] (0,0) -- (0,3.1) node[above, font=\footnotesize] {canary traffic \%};
  \foreach \y/\p in {0.55/5, 1.35/25, 2.15/50, 2.85/100}
    \draw[packtgray!40, dashed] (0,\y) -- (13.0,\y) node[pos=0, left, note] {\p\%};
  % traffic staircase (canary share)
  \draw[very thick, packtorange] (0.2,0.15) -- (1.2,0.15) -- (1.2,0.55) -- (3.0,0.55)
    -- (3.0,1.35) -- (5.2,1.35) -- (5.2,2.15) -- (7.4,2.15) -- (7.4,2.85) -- (11.4,2.85);
  \node[modcap, anchor=south west] at (0.2,3.0) {traffic staircase with metric gates};
  % analysis windows
  \foreach \x/\ok in {1.2/1, 3.0/1, 5.2/1, 7.4/1} {
    \draw[<->, thick, blue!60] (\x-0.15,0.30) -- (\x+1.55,0.30);
  }
  \node[note, align=center] at (2.1,0.02) {gate 1: SPRT\\passes};
  \node[note, align=center] at (4.1,0.02) {gate 2: p95 within\\1.25x baseline};
  \node[note, align=center] at (6.3,0.02) {gate 3: error delta\\$< \delta$};
  \node[note, align=center] at (9.4,0.02) {gate 4: full-window\\analysis, then 100\%};
  % rollback example branch
  \draw[very thick, red!70, dashed] (3.0,1.35) -- (3.0,0.15) -- (4.4,0.15);
  \node[note, text=red!70, align=left] at (4.6,0.35) {if any gate fails: LLR $\ge \ln\frac{1-\beta}{\alpha}$\\$\Rightarrow$ automated rollback to 0\%};
  % completion marker
  \node[flowstep, fill=green!10] at (12.2,2.85) {promoted};
\end{tikzpicture}
\caption{Canary rollout timeline. Each step holds for a fixed analysis
window while the gate (error-rate SPRT plus latency guardrails,
Listing~\ref{lst:ch17-canary}) evaluates; evidence crossing the upper
boundary rolls back automatically, the lower boundary promotes.}
\label{fig:ch17-canary}
\end{figure}

\subsection{Edge \& TLS: Where Bytes Get Terminated}
\label{sec:ch17-edge}

\textbf{TLS termination placement} is a three-way decision with real
consequences, not a default. \emph{Edge termination} (decrypt at
CDN/LB/gateway, plaintext inside) minimizes certificate sprawl and
enables L7 policy, but leaves the last hop plaintext --- acceptable only
on networks you genuinely control and increasingly non-compliant.
\emph{Re-encryption} (terminate at the edge, open a \emph{second} TLS
session to the backend) keeps policy visibility \emph{and} encrypted
transit at the price of two certificate lifecycles. \emph{Passthrough}
(L4 forwarding of the encrypted stream to the service) gives end-to-end
encryption with zero edge visibility --- no WAF, no L7 routing, no JWT
offload. The common production answer for APIs: terminate at the edge
with automated public certificates, re-encrypt to services with an
internal CA (a mesh gives you this as mTLS ``for free''), and reserve
passthrough for regulated payloads. \textbf{Certificate automation} is
non-negotiable either way: ACME (Let's Encrypt) at the edge,
cert-manager or mesh-issued certificates inside --- a manually renewed
certificate is an outage with a calendar date
(Chapter~\ref{ch:api_security}).

\textbf{CDN behavior for APIs} is mostly about what \emph{not} to cache.
Cacheable: truly anonymous, truly idempotent GETs with explicit
\texttt{Cache-Control} (public reference data, OpenAPI documents, JWKS).
Never cacheable: authenticated responses (unless the cache key includes
the auth scope --- rarely worth it), anything with \texttt{Set-Cookie},
and personalized content. The \textbf{cache key} discipline: default
keys ignore query strings and headers, so filtered list endpoints
(Chapter~\ref{ch:pagination}) must either opt out or list their
query/header inputs in the key; \texttt{Vary} mishandling is the classic
tenant-data-leak vector. \textbf{Purge} semantics matter more than TTLs
for mutable resources: if you cannot purge by key or tag on write, keep
TTLs in seconds and accept staleness as a documented contract
(Chapter~\ref{ch:idempotency}'s caching half). \textbf{WAF essentials}:
managed rule sets (OWASP core rules) in front of the gateway, tuned to
block mode gradually (detect first --- API payloads trip SQLi rules on
legitimate JSON), plus body-size limits and request-rate anomaly rules
as the coarse outer net before the gateway's precise per-key limiting
(Chapter~\ref{ch:rate_limiting}).

\begin{definition}[North-south vs.\ east-west traffic]
\emph{North-south} traffic crosses the platform boundary: client to
edge, edge to gateway, gateway to first service. \emph{East-west}
traffic stays inside: service to service, service to database. The
distinction selects the control point: gateways, WAFs, and CDNs see only
north-south; meshes, internal firewalls, and mTLS see east-west. A
policy applied at the wrong tier (e.g., per-client rate limits enforced
only east-west) protects nothing, because the abuse never reaches it.
\end{definition}

\begin{figure}[h]
\centering
\begin{tikzpicture}[font=\small, node distance=4mm]
  \node[flowstep, fill=packtgray!8] (client) {API consumer\\(mobile / partner)};
  \node[flowstep, fill=packtorange!12, right=14mm of client] (cdn) {Provider edge\\CDN + WAF\\TLS terminate (public CA)};
  \node[flowstep, fill=blue!8, right=12mm of cdn] (gw) {API gateway\\(Kong/Envoy)\\JWT check, rate limit, canary split};
  \node[flowstep, fill=green!8, right=12mm of gw] (mesh) {mesh sidecar\\mTLS re-encrypt\\retries, telemetry};
  \node[flowstep, fill=packtgray!8, right=12mm of mesh] (svc) {Orders API pod\\FastAPI/uvicorn\\non-root, read-only fs};
  \node[flowstep, fill=green!8, below=10mm of svc] (peer) {Inventory API pod};
  \draw[->, thick] (client) -- node[above, note]{HTTPS 443} (cdn);
  \draw[->, thick] (cdn) -- node[above, note]{HTTPS, WAF-clean} (gw);
  \draw[->, thick] (gw) -- node[above, note]{re-encrypted} (mesh);
  \draw[->, thick] (mesh) -- node[above, note]{mTLS} (svc);
  \draw[->, thick] (svc) -- node[right, note]{east-west} (peer);
  % boundary boxes (drawn last; unfilled dashed borders only)
  \node[draw=packtorange, dashed, fit=(cdn), inner sep=6pt, label={[note]south:provider edge}] {};
  \node[draw=blue!60, dashed, fit=(gw)(mesh)(svc)(peer), inner sep=8pt, label={[note]south east:your platform}] {};
  % north-south brace
  \draw[decorate, decoration={brace, amplitude=5pt}, packtgray]
    ([yshift=16mm]client.west) -- node[above=6pt, note]{north-south} ([yshift=16mm]gw.east);
\end{tikzpicture}
\caption{Request path through edge, gateway, and mesh. TLS terminates
at the provider edge (public certificate, WAF inspection), is
re-encrypted to the gateway, and becomes mTLS east-west; each box owns
exactly the policies its position can see.}
\label{fig:ch17-requestpath}
\end{figure}

%% ----------------------------------------------------------------------------
\section{Production-Ready Code Implementation}
\label{sec:ch17-code}

This section contains every artifact the chapter discusses: the
Dockerfile and its build-context guard, the full Kubernetes manifest set
(reviewed line by line, never applied), and the three runnable,
offline, deterministic programs --- the graceful-shutdown proof, the
probe-semantics simulator, and the canary analyzer.

\subsection{Listing 17.1 --- Production Multi-Stage Dockerfile}
\label{sec:ch17-dockerfile}
\label{lst:ch17-dockerfile}

\begin{minted}{text}
# ---------- Stage 1: build (discarded) ----------
FROM python:3.12-slim AS build
ENV PIP_NO_CACHE_DIR=1 \
    PIP_DISABLE_PIP_VERSION_CHECK=1
WORKDIR /build
RUN python -m venv /opt/venv
ENV PATH="/opt/venv/bin:$PATH"
COPY requirements.txt .
RUN pip install -r requirements.txt

# ---------- Stage 2: runtime (shipped) ----------
# Pin by digest for real releases; resolve it with:
#   docker buildx imagetools inspect python:3.12-slim
FROM python:3.12-slim AS runtime
ENV PYTHONDONTWRITEBYTECODE=1 \
    PYTHONUNBUFFERED=1 \
    PATH="/opt/venv/bin:$PATH"
RUN groupadd --system app \
 && useradd --system --gid app --uid 10001 --no-create-home app
WORKDIR /app
COPY --from=build /opt/venv /opt/venv
COPY app/ ./app/
USER app
EXPOSE 8080
HEALTHCHECK --interval=30s --timeout=3s --start-period=15s --retries=3 \
  CMD ["python", "-c", "import urllib.request,sys; \
       sys.exit(0 if urllib.request.urlopen('http://127.0.0.1:8080/healthz', timeout=2).status == 200 else 1)"]
ENTRYPOINT ["python", "-m", "uvicorn", "app.main:app", \
            "--host", "0.0.0.0", "--port", "8080"]
\end{minted}

Line-by-line, with the reason each line earns its place:

\begin{itemize}
  \item \textbf{Stage split} (\texttt{AS build} / \texttt{AS runtime}):
        the build stage owns \texttt{pip} and the resolver; only the
        finished virtual environment crosses into the runtime stage via
        \texttt{COPY --from=build}. Compilers, caches, and package
        managers never ship.
  \item \texttt{PIP\_NO\_CACHE\_DIR=1}: the pip cache is pure dead weight
        in a layer --- every cached wheel doubles the cost of the layer
        that installed it.
  \item \textbf{venv at a fixed path} (\texttt{/opt/venv}): a virtual
        environment is a relocatable, root-owned-by-build artifact; the
        runtime stage sets \texttt{PATH} once and never invokes pip
        again. \texttt{PYTHONDONTWRITEBYTECODE} plus a read-only root
        filesystem agree: nothing writes \texttt{.pyc} at runtime.
  \item \textbf{Cache ordering}: \texttt{COPY requirements.txt} and the
        \texttt{pip install} precede \texttt{COPY app/} so that a
        source-only change reuses the entire dependency layer; a
        dependency change correctly invalidates from the manifest copy
        downward.
  \item \textbf{Digest pinning} (commented): tags are mutable registry
        pointers; a digest is a content hash. Production pipelines
        resolve the digest at build time (\texttt{imagetools inspect})
        and record it, so ``what we tested'' and ``what is running'' are
        the same manifest.
  \item \textbf{Non-root}: a dedicated system user (fixed UID 10001 so
        Kubernetes \texttt{runAsUser} can assert the same number) is the
        last instruction before \texttt{ENTRYPOINT}; everything after
        \texttt{USER app} runs unprivileged.
  \item \textbf{\texttt{HEALTHCHECK}}: the image carries its own health
        opinion --- stdlib \texttt{urllib} against \texttt{/healthz},
        no \texttt{curl} dependency (slim images lack it; installing it
        to run a health check would grow the attack surface to check the
        attack surface). Kubernetes ignores \texttt{HEALTHCHECK} in favor
        of probes; Docker and compose honor it.
  \item \textbf{Exec-form \texttt{ENTRYPOINT}}: the JSON array form means
        uvicorn \emph{is} PID~1 --- no shell wrapper to swallow
        \texttt{SIGTERM} (Section~\ref{sec:ch17-process}). Binding
        \texttt{0.0.0.0} is correct \emph{inside} a container: the
        network namespace is already the isolation boundary.
\end{itemize}

The build-context guard, without which the cache ordering above is
fiction (a stray log file changing would invalidate the source layer) and
secrets can leak into layers:

\begin{minted}{text}
# .dockerignore
.git
.venv
__pycache__/
*.pyc
.pytest_cache/
tests/
datasets/
*.md
.env
.env.*
Dockerfile
.dockerignore
\end{minted}

Runtime hardening completes the artifact: \texttt{docker run --read-only
--tmpfs /tmp:rw,noexec,nosuid,size=16m \ldots} gives the process a
read-only root filesystem with a tiny, non-executable scratch space; the
Kubernetes equivalent (\texttt{readOnlyRootFilesystem: true} plus an
\texttt{emptyDir} mounted at \texttt{/tmp}) appears in
Listing~\ref{lst:ch17-deployment}.

\begin{examplebox}
\textbf{Size as a report card.} The same Orders API built three ways on
the authors' reference machine: full \texttt{python:3.12} base:
$\sim$1.02~GB, 3 high-severity and 41 lower findings from a standard
scanner. \texttt{python:3.12-slim} multi-stage (Listing 17.1):
$\sim$165~MB, 0 high, 9 low. Distroless runtime: $\sim$78~MB, 0 high,
2 low --- and no shell, which flips several scanner findings from
``exploitable'' to ``unreachable.'' Image size is not vanity; it is
transfer time, cold-start time, and vulnerability count in one number.
\end{examplebox}

\subsection{Listing 17.2 --- Kubernetes Manifest Set (Reviewed, Not Applied)}
\label{sec:ch17-manifests}

The set below deploys the Orders API with a stable track and a canary
track. Read each block with the justifications; the design review in
Section~\ref{sec:ch17-lab} walks the same fields as a meeting exercise.

\subsubsection*{Deployment (stable track)}
\label{lst:ch17-deployment}

\begin{minted}{yaml}
apiVersion: apps/v1
kind: Deployment
metadata:
  name: orders-api
  labels: {app: orders-api, track: stable}
spec:
  replicas: 4
  revisionHistoryLimit: 5
  progressDeadlineSeconds: 600
  strategy:
    type: RollingUpdate
    rollingUpdate:
      maxSurge: 1
      maxUnavailable: 0
  selector:
    matchLabels: {app: orders-api, track: stable}
  template:
    metadata:
      labels: {app: orders-api, track: stable, version: "2.4.1"}
    spec:
      terminationGracePeriodSeconds: 45
      securityContext:
        runAsNonRoot: true
        runAsUser: 10001
        fsGroup: 10001
      containers:
        - name: api
          image: registry.northwind-robotics.example/orders-api:2.4.1
          ports:
            - {name: http, containerPort: 8080}
          env:
            - {name: UVICORN_WORKERS, value: "2"}
            - {name: UVICORN_TIMEOUT_GRACEFUL_SHUTDOWN, value: "30"}
            - name: DATABASE_URL
              valueFrom:
                secretKeyRef: {name: orders-api-secrets, key: database-url}
          resources:
            requests: {cpu: 250m, memory: 256Mi}
            limits: {memory: 256Mi}     # no CPU limit: throttle pitfall
          securityContext:
            allowPrivilegeEscalation: false
            readOnlyRootFilesystem: true
            capabilities: {drop: ["ALL"]}
          volumeMounts:
            - {name: tmp, mountPath: /tmp}
          lifecycle:
            preStop:
              sleep: {seconds: 8}       # cover endpoint-removal lag
          startupProbe:
            httpGet: {path: /healthz/startup, port: http}
            periodSeconds: 5
            failureThreshold: 12        # budget: 60 s to boot
          readinessProbe:
            httpGet: {path: /healthz/ready, port: http}
            periodSeconds: 5
            failureThreshold: 3
          livenessProbe:
            httpGet: {path: /healthz/live, port: http}
            periodSeconds: 10
            failureThreshold: 3
      volumes:
        - name: tmp
          emptyDir: {}
\end{minted}

Field by field: \texttt{replicas: 4} pairs with \texttt{maxSurge: 1,
maxUnavailable: 0} --- capacity never dips below 4 during a rollout, at
the price of one surplus pod; the arithmetic for other sizes is in
Table~\ref{tab:ch17-rollout}. \texttt{progressDeadlineSeconds: 600} makes
a wedged rollout (e.g., a readiness path renamed in the new version)
\emph{fail loudly} after ten minutes instead of hanging silently.
\texttt{revisionHistoryLimit: 5} keeps five old ReplicaSets for instant
rollback without unbounded clutter. \texttt{terminationGracePeriodSeconds:
45} is the shutdown budget: 8~s preStop sleep $+$ 30~s uvicorn graceful
timeout $+$ margin; the two values must be set together --- a 30~s
default grace against a 45~s drain need is a dropped-request generator.
The pod-level \texttt{securityContext} pins the UID the image was built
with (mutual assertion: Kubernetes refuses a container whose image tries
to run as root). \texttt{resources}: memory request equals limit (memory
cannot be throttled; a mismatch is pure OOM risk), CPU has a request
(scheduling and CFS shares) but \emph{no limit} --- the throttling
pitfall of Section~\ref{sec:ch17-k8s}. \texttt{DATABASE\_URL} comes from
a Secret reference; nothing secret is ever a literal in a manifest
(Chapter~\ref{ch:api_security}). \texttt{readOnlyRootFilesystem: true}
plus the \texttt{emptyDir} at \texttt{/tmp} is the runtime-fleet version
of the Dockerfile's discipline. The \texttt{preStop} \texttt{sleep} uses
the native sleep action (Kubernetes $\ge$ 1.30); on older clusters the
equivalent is \texttt{exec: \{command: ["sh", "-c", "sleep 8"]\}}, which
requires a shell in the image --- one concrete reason to choose slim over
distroless. The three probes implement Table~\ref{tab:ch17-probes}
verbatim: startup gates the others with a 60-second boot budget;
readiness is the only probe allowed to remove endpoints; liveness checks
a dependency-free path so only process-local death can restart the pod.

\subsubsection*{Service}
\label{lst:ch17-service}

\begin{minted}{yaml}
apiVersion: v1
kind: Service
metadata:
  name: orders-api
spec:
  selector:
    app: orders-api        # no 'track' label: spans stable AND canary
  ports:
    - {name: http, port: 80, targetPort: http}
\end{minted}

The selector's omission is the canary mechanism: the Service load
balances across both tracks in proportion to ready pod count, so one
canary pod against four stable pods receives $\approx 20\%$ of traffic.
Pod-count splitting is the coarse, zero-infrastructure option; arbitrary
weights (and header/tenant-based pinning) require gateway weighted
clusters --- Envoy's first-class feature in
Table~\ref{tab:ch17-gateway-matrix}.

\subsubsection*{HorizontalPodAutoscaler}
\label{lst:ch17-hpa}

\begin{minted}{yaml}
apiVersion: autoscaling/v2
kind: HorizontalPodAutoscaler
metadata:
  name: orders-api
spec:
  scaleTargetRef: {apiVersion: apps/v1, kind: Deployment, name: orders-api}
  minReplicas: 4
  maxReplicas: 20
  metrics:
    - type: Resource
      resource:
        name: cpu
        target: {type: Utilization, averageUtilization: 70}
  behavior:
    scaleUp:
      stabilizationWindowSeconds: 0
      policies:
        - {type: Percent, value: 100, periodSeconds: 60}
    scaleDown:
      stabilizationWindowSeconds: 300
      policies:
        - {type: Pods, value: 1, periodSeconds: 60}
\end{minted}

Utilization targets CPU \emph{request} (70\% of 250m), which is why the
request must reflect reality. The asymmetric behavior is the point:
scale up instantly (a flash crowd does not wait five minutes) but scale
down through a 300-second stabilization window and at most one pod per
minute, because every scale-down is a termination event with in-flight
requests attached --- the HPA's politeness and the preStop hook are two
halves of the same courtesy. \texttt{minReplicas: 4} interacts with the
PDB below: voluntary disruptions always leave three.

\subsubsection*{PodDisruptionBudget}
\label{lst:ch17-pdb}

\begin{minted}{yaml}
apiVersion: policy/v1
kind: PodDisruptionBudget
metadata:
  name: orders-api
spec:
  minAvailable: 3
  selector:
    matchLabels: {app: orders-api}
\end{minted}

\texttt{minAvailable: 3} means node drains and cluster upgrades must
leave three pods serving; an eviction that would violate the budget
\emph{blocks} (the drain waits, it does not force). The classic
misconfiguration is \texttt{minAvailable} equal to the replica count,
which deadlocks every drain; the second classic is omitting the PDB
entirely, which lets \texttt{kubectl drain} evict the whole fleet of a
single-replica API during a kernel upgrade.

\subsubsection*{Canary Deployment (variant)}
\label{lst:ch17-canary-deploy}

\begin{minted}{yaml}
apiVersion: apps/v1
kind: Deployment
metadata:
  name: orders-api-canary
spec:
  replicas: 1                      # 1 canary vs 4 stable = ~20% traffic
  strategy:
    type: RollingUpdate
    rollingUpdate: {maxSurge: 1, maxUnavailable: 0}
  selector:
    matchLabels: {app: orders-api, track: canary}
  template:
    metadata:
      labels: {app: orders-api, track: canary, version: "2.5.0-rc.1"}
    spec:
      # identical pod spec to the stable Deployment except:
      containers:
        - name: api
          image: registry.northwind-robotics.example/orders-api:2.5.0-rc.1
          # ... probes, resources, securityContext, lifecycle unchanged ...
\end{minted}

The canary Deployment differs in exactly three places: name, replica
count, and the \texttt{track}/\texttt{version} labels that keep the two
ReplicaSets distinct and let observability slice metrics by version
(Chapter~\ref{ch:observability}). Promotion is \texttt{kubectl scale}
steps on the two Deployments (2/3, 3/2, 4/1) gated by
Listing~\ref{lst:ch17-canary}'s analyzer; rollback is scaling the canary
to zero --- fast precisely because the expand-and-contract schema
discipline of Chapter~\ref{ch:versioning} guaranteed the old binary still
understands the data.

\subsection{Listing 17.3 --- The Server Under Test (\texttt{orders\_server.py})}
\label{sec:ch17-server}
\label{lst:ch17-server}

The FastAPI Orders API in its deployment-relevant form: in-flight request
tracking via middleware (the ledger that lets us \emph{prove} a drain),
a fast write path, a deliberately slow five-second report endpoint to be
caught mid-flight, and a probe endpoint~\cite{fastapidocs}.

\begin{minted}{python}
"""Northwind Robotics Orders API -- used to demonstrate graceful shutdown.

Run standalone (``python orders_server.py``) or under the driver in
``shutdown_demo.py``. Every response carries three diagnostic headers so an
outside observer can prove what happened to in-flight work during a
termination signal:

* ``X-Request-Id`` -- monotonically increasing request identifier.
* ``X-In-Flight-Remaining`` -- requests still in flight after this one.
* ``X-Elapsed-Ms`` -- server-side wall time for the request.
"""

from __future__ import annotations

import asyncio
import itertools
import time

from fastapi import FastAPI, Request, Response

app = FastAPI(title="Northwind Robotics Orders API")

IN_FLIGHT: dict[int, str] = {}
_request_ids = itertools.count(1)


@app.middleware("http")
async def track_in_flight(request: Request, call_next) -> Response:
    """Ledger middleware: register on entry, deregister on exit."""
    rid = next(_request_ids)
    IN_FLIGHT[rid] = f"{request.method} {request.url.path}"
    started = time.perf_counter()
    try:
        response = await call_next(request)
    finally:
        IN_FLIGHT.pop(rid, None)
    response.headers["X-Request-Id"] = str(rid)
    response.headers["X-In-Flight-Remaining"] = str(len(IN_FLIGHT))
    response.headers["X-Elapsed-Ms"] = f"{(time.perf_counter() - started) * 1000:.0f}"
    return response


@app.get("/healthz")
async def healthz() -> dict[str, bool]:
    """Cheap probe endpoint: process is alive and the event loop turns."""
    return {"ok": True}


@app.post("/orders", status_code=201)
async def create_order() -> dict[str, str]:
    """Fast path: a normal order write."""
    await asyncio.sleep(0.05)  # simulated database round trip
    return {"order_id": "ord-1042", "status": "accepted"}


@app.get("/orders/export")
async def export_orders() -> dict[str, object]:
    """Slow path: a five-second report, one simulated chunk per second.

    This endpoint exists to be *in flight* when the termination signal
    arrives; whether it completes is the whole point of the demo.
    """
    chunks = 0
    for _ in range(5):
        await asyncio.sleep(1.0)  # one chunk of report work per second
        chunks += 1
    return {
        "report": "orders",
        "chunks_completed": chunks,
        "in_flight_at_finish": len(IN_FLIGHT),
    }


if __name__ == "__main__":
    import uvicorn

    uvicorn.run(app, host="127.0.0.1", port=8571, log_level="warning")
\end{minted}

\subsection{Listing 17.4 --- The Graceful-Shutdown Proof (\texttt{shutdown\_demo.py})}
\label{sec:ch17-shutdown}
\label{lst:ch17-shutdown}

The driver starts the server as a child process, puts the five-second
export in flight, sends the termination signal, and then verifies the
four-act protocol of Section~\ref{sec:ch17-process} from the
\emph{client's} side: a new request must be refused (stop accepting), the
in-flight request must complete with HTTP~200 (drain), and the process
must exit on its own. On Linux the signal is \texttt{SIGTERM}; Windows
has no POSIX signal delivery between processes, so the driver sends
\texttt{CTRL\_BREAK\_EVENT} to the child's console process group ---
uvicorn installs a \texttt{SIGBREAK} handler on Windows that enters
exactly the same graceful-exit path \texttt{SIGTERM} triggers on Linux,
which is the path that matters, because the production artifact is a
Linux container receiving \texttt{SIGTERM} from the kubelet.

\begin{minted}{python}
"""Prove uvicorn's drain semantics: terminate a server mid-request.

Offline and deterministic in behavior: the driver starts
``orders_server.py`` as a child process, launches a slow ``/orders/export``
request (about five seconds), sends the termination signal roughly one
second in, and then verifies the three properties a graceful shutdown must
have, in order:

1. STOP ACCEPTING -- a brand-new request after the signal is refused.
2. DRAIN          -- the in-flight export finishes with HTTP 200.
3. EXIT           -- the server process terminates on its own.

Usage (PowerShell):  python shutdown_demo.py
Exit code 0 means all three properties held; 1 means a drain violation.
"""

from __future__ import annotations

import json
import os
import signal
import subprocess
import sys
import threading
import time
import urllib.error
import urllib.request
from pathlib import Path

HOST, PORT = "127.0.0.1", 8571
BASE = f"http://{HOST}:{PORT}"
SERVER_SCRIPT = Path(__file__).resolve().parent / "orders_server.py"

_step = 0


def log(message: str) -> None:
    """Print a numbered timeline line; numbering makes ordering explicit."""
    global _step
    _step += 1
    print(f"[{_step:02d}] {message}")


def get(path: str) -> tuple[int, dict[str, str], dict]:
    """Plain-stdlib GET returning (status, headers, parsed JSON body).

    Header keys are normalized to lowercase because ASGI servers emit
    lowercase field names on the wire.
    """
    with urllib.request.urlopen(f"{BASE}{path}", timeout=25.0) as resp:
        headers = {k.lower(): v for k, v in resp.headers.items()}
        return resp.status, headers, json.loads(resp.read())


def wait_ready(server: subprocess.Popen, timeout: float = 15.0) -> None:
    """Poll /healthz until the server answers or dies trying."""
    deadline = time.monotonic() + timeout
    while time.monotonic() < deadline:
        if server.poll() is not None:
            raise RuntimeError(f"server exited early, code {server.returncode}")
        try:
            status, _, _ = get("/healthz")
            if status == 200:
                return
        except (urllib.error.URLError, ConnectionError, OSError):
            time.sleep(0.25)
    raise TimeoutError("server did not become ready in time")


def main() -> int:
    creationflags = 0
    if os.name == "nt":
        # Give the child its own console process group so CTRL_BREAK_EVENT
        # reaches it (and only it) on Windows.
        creationflags = subprocess.CREATE_NEW_PROCESS_GROUP
    server = subprocess.Popen(
        [sys.executable, str(SERVER_SCRIPT)],
        stdout=subprocess.DEVNULL,
        stderr=subprocess.DEVNULL,
        creationflags=creationflags,
    )
    try:
        wait_ready(server)
        log("server ready; /healthz answers 200")

        outcome: dict[str, object] = {}

        def slow_call() -> None:
            try:
                status, headers, body = get("/orders/export")
                outcome["status"] = status
                outcome["elapsed_ms"] = headers.get("x-elapsed-ms", "?")
                outcome["chunks"] = body.get("chunks_completed")
            except (urllib.error.URLError, ConnectionError, OSError) as exc:
                outcome["error"] = f"{type(exc).__name__}: {exc}"

        worker = threading.Thread(target=slow_call, daemon=True)
        worker.start()
        time.sleep(1.2)  # the five-second export is now mid-flight

        if os.name == "nt":
            sig_name, sig = "CTRL_BREAK_EVENT", signal.CTRL_BREAK_EVENT
        else:
            sig_name, sig = "SIGTERM", signal.SIGTERM
        log(f"sending {sig_name} while the export is in flight")
        os.kill(server.pid, sig)

        time.sleep(0.3)
        log("attempting a NEW request after the signal")
        try:
            get("/healthz")
            new_blocked = False
            log("UNEXPECTED: a new request was accepted after the signal")
        except (urllib.error.URLError, ConnectionError, OSError) as exc:
            new_blocked = True
            log(f"new request refused as expected ({type(exc).__name__})")

        worker.join(timeout=25.0)
        if "error" in outcome:
            log(f"in-flight export FAILED: {outcome['error']}")
        else:
            log(
                "in-flight export completed: "
                f"HTTP {outcome.get('status')}, "
                f"elapsed {outcome.get('elapsed_ms')} ms, "
                f"chunks {outcome.get('chunks')}/5"
            )

        returncode = server.wait(timeout=25.0)
        log(f"server process exited on its own (code {returncode})")

        drained = outcome.get("status") == 200 and outcome.get("chunks") == 5
        print()
        print("VERDICT")
        print(f"  1. stop accepting first : {'PASS' if new_blocked else 'FAIL'}")
        print(f"  2. drain in-flight work : {'PASS' if drained else 'FAIL'}")
        print(f"  3. clean process exit   : {'PASS' if returncode is not None else 'FAIL'}")
        return 0 if (new_blocked and drained) else 1
    finally:
        if server.poll() is None:  # never leave an orphaned server behind
            server.kill()
            server.wait(timeout=10)


if __name__ == "__main__":
    raise SystemExit(main())
\end{minted}

Observed output (Windows 11, Python 3.13, uvicorn 0.49; on Linux the
signal name reads \texttt{SIGTERM} and the exit code is 0):

\begin{minted}{text}
[01] server ready; /healthz answers 200
[02] sending CTRL_BREAK_EVENT while the export is in flight
[03] attempting a NEW request after the signal
[04] new request refused as expected (URLError)
[05] in-flight export completed: HTTP 200, elapsed 5036 ms, chunks 5/5
[06] server process exited on its own (code 3)

VERDICT
  1. stop accepting first : PASS
  2. drain in-flight work : PASS
  3. clean process exit   : PASS
\end{minted}

Read line [05] carefully: the export was signaled at roughly 1.2~seconds
into a five-second job and still reports all five chunks and an elapsed
time of 5036~ms --- the drain completed \emph{after} the refusal of new
work, which is the only ordering that is both safe and useful. The exit
code (3 on Windows, 0 on Linux) is platform plumbing; the contract is
that the process exits \emph{itself}, before the supervisor's patience
(\texttt{terminationGracePeriodSeconds}) runs out.

\subsection{Listing 17.5 --- Probe-Semantics Simulator (\texttt{probe\_simulator.py})}
\label{sec:ch17-probesim}
\label{lst:ch17-probesim}

Table~\ref{tab:ch17-probes} stated the rules; this simulator enforces
them. It models one Pod under the kubelet's probe logic over a virtual
one-second-tick clock and a scripted fault timeline (slow boot, then a
35-second dependency outage, then a transient deadlock), and prints the
two things a caller can observe: endpoint membership (does this Pod get
traffic?) and restarts. Scenario~B removes the startup probe --- the
liveness-as-boot-check anti-pattern --- and demonstrates
CrashLoopBackOff on the same world.

\begin{minted}{python}
"""Kubernetes probe-semantics simulator (offline, stdlib-only, deterministic).

Two scenarios run over the same world:

  A. WITH a startup probe: slow boots are tolerated; a deadlock is
     detected by liveness and fixed by a restart; a dependency outage
     only removes the Pod from Service endpoints (no restart).
  B. WITHOUT a startup probe (the copy-paste anti-pattern where the
     liveness probe doubles as the startup check): a 45-second cold boot
     exceeds the liveness budget, so the kubelet kills the container
     before it ever serves -- CrashLoopBackOff.

The clock is virtual (one-second ticks), so every run prints byte-for-byte
identical output.
"""

from __future__ import annotations

from dataclasses import dataclass, field
from enum import Enum

SIM_SECONDS = 240
MAX_RESTARTS = 2  # afterwards the kubelet backs off (CrashLoopBackOff)


class Mode(Enum):
    BOOTING = "booting"          # process alive, app still initializing
    HEALTHY = "healthy"
    DEP_DOWN = "dep_down"        # dependency down: /ready fails, /live fine
    DEADLOCKED = "deadlocked"    # event loop stuck: every probe times out


@dataclass(frozen=True)
class ProbeSpec:
    name: str
    period_seconds: int
    failure_threshold: int
    success_threshold: int = 1


STARTUP = ProbeSpec("startup", period_seconds=5, failure_threshold=12)  # 60 s
READINESS = ProbeSpec("readiness", period_seconds=5, failure_threshold=3)
LIVENESS = ProbeSpec("liveness", period_seconds=10, failure_threshold=3)

# The fault environment: (start_s, end_s, mode once booted).
TIMELINE: list[tuple[int, int, Mode]] = [
    (0, 60, Mode.HEALTHY),
    (60, 95, Mode.DEP_DOWN),       # 35 s database outage
    (95, 150, Mode.HEALTHY),
    (150, 10_000, Mode.DEADLOCKED),  # transient deadlock; a restart cures it
]


def timeline_mode(t: int) -> Mode:
    for start, end, mode in TIMELINE:
        if start <= t < end:
            return mode
    return Mode.HEALTHY


@dataclass
class SimResult:
    label: str
    events: list[str] = field(default_factory=list)
    restarts: int = 0
    traffic_seconds: int = 0
    ready_at_end: bool = False


def run_scenario(label: str, use_startup_probe: bool, boot_seconds: int,
                 reboot_seconds: int) -> SimResult:
    """Simulate one Pod. Returns an event log plus routing statistics."""
    result = SimResult(label=label)

    boot_started = 0
    boot_len = boot_seconds
    booted = False          # startup check has passed at least once
    ready = False           # Pod is a member of the Service endpoints
    deadlock_cleared = False  # restarts cure process-local deadlocks
    crashlooped = False

    startup_fail = ready_fail = ready_ok = live_fail = 0

    def mode_at(t: int) -> Mode:
        if t - boot_started < boot_len:
            return Mode.BOOTING
        mode = timeline_mode(t)
        if mode is Mode.DEADLOCKED and deadlock_cleared:
            return Mode.HEALTHY
        return mode

    def restart(t: int, why: str) -> None:
        nonlocal boot_started, boot_len, booted, ready
        nonlocal startup_fail, ready_fail, ready_ok, live_fail
        nonlocal deadlock_cleared, crashlooped
        result.restarts += 1
        if result.restarts > MAX_RESTARTS:
            crashlooped = True
            result.events.append(
                f"t={t:3d}s  restart budget exhausted -> CrashLoopBackOff; "
                "kubelet backs off, Pod never becomes ready"
            )
            return
        result.events.append(f"t={t:3d}s  RESTART #{result.restarts} ({why})")
        boot_started, boot_len = t, reboot_seconds
        booted = ready = False
        startup_fail = ready_fail = ready_ok = live_fail = 0
        deadlock_cleared = True

    def set_ready(t: int, value: bool, why: str) -> None:
        nonlocal ready
        if ready != value:
            ready = value
            action = "ADDED TO" if value else "REMOVED FROM"
            result.events.append(
                f"t={t:3d}s  endpoints: Pod {action} Service rotation ({why})"
            )

    for t in range(1, SIM_SECONDS + 1):
        if crashlooped:
            break
        mode = mode_at(t)

        # --- startup probe (gates the other two) -------------------------
        if use_startup_probe and not booted and t % STARTUP.period_seconds == 0:
            if mode is not Mode.BOOTING:
                booted = True
                startup_fail = 0
                result.events.append(
                    f"t={t:3d}s  startup probe: first success -> "
                    "liveness/readiness probes begin"
                )
            else:
                startup_fail += 1
                if startup_fail >= STARTUP.failure_threshold:
                    restart(t, "startup probe exhausted its budget")

        probes_active = booted if use_startup_probe else True

        # --- readiness probe: endpoint membership only --------------------
        if probes_active and t % READINESS.period_seconds == 0:
            if mode is Mode.HEALTHY:
                ready_fail = 0
                ready_ok += 1
                if ready_ok >= READINESS.success_threshold:
                    set_ready(t, True, "readiness success")
            else:
                ready_ok = 0
                ready_fail += 1
                if ready_fail >= READINESS.failure_threshold:
                    set_ready(t, False, "readiness failing -- no restart!")

        # --- liveness probe: restart authority ----------------------------
        if probes_active and t % LIVENESS.period_seconds == 0:
            if mode in (Mode.HEALTHY, Mode.DEP_DOWN):
                live_fail = 0
            else:  # BOOTING (no startup probe) or DEADLOCKED: timeout
                live_fail += 1
                result.events.append(
                    f"t={t:3d}s  liveness probe: failure {live_fail}/"
                    f"{LIVENESS.failure_threshold} (mode={mode.value})"
                )
                if live_fail >= LIVENESS.failure_threshold:
                    why = ("deadlock detected" if mode is Mode.DEADLOCKED
                           else "killed while still booting")
                    restart(t, why)

        if ready:
            result.traffic_seconds += 1

    result.ready_at_end = ready
    return result


def report(result: SimResult) -> None:
    print("=" * 72)
    print(f"SCENARIO {result.label}")
    print("=" * 72)
    for line in result.events:
        print(" ", line)
    print("-" * 72)
    print(f"  restarts                    : {result.restarts}")
    print(f"  seconds receiving traffic   : {result.traffic_seconds} / {SIM_SECONDS}")
    print(f"  ready at end of simulation  : {result.ready_at_end}")
    print()


def main() -> None:
    print("Kubernetes probe-semantics simulator")
    print(f"probes: startup {STARTUP.period_seconds}s x{STARTUP.failure_threshold}, "
          f"readiness {READINESS.period_seconds}s x{READINESS.failure_threshold}, "
          f"liveness {LIVENESS.period_seconds}s x{LIVENESS.failure_threshold}")
    print("faults: 18s boot | dep outage t=60..95 | deadlock at t=150 "
          "(cured by restart)")
    print()

    with_startup = run_scenario(
        "A -- WITH startup probe (boot 18s tolerated)",
        use_startup_probe=True, boot_seconds=18, reboot_seconds=12,
    )
    report(with_startup)

    without_startup = run_scenario(
        "B -- WITHOUT startup probe (liveness doubles as boot check, boot 45s)",
        use_startup_probe=False, boot_seconds=45, reboot_seconds=45,
    )
    report(without_startup)

    print("LESSONS")
    print("  * DEP_DOWN caused readiness removal but ZERO restarts --")
    print("    restarting cannot fix someone else's outage.")
    print("  * DEADLOCKED caused a restart after 3 consecutive liveness")
    print("    failures -- liveness exists for process-local death.")
    print("  * Without a startup probe, a slow boot is indistinguishable from")
    print("    a deadlock, and the liveness budget becomes a boot deadline.")


if __name__ == "__main__":
    main()
\end{minted}

Observed output (identical on every run --- the clock is virtual):

\begin{minted}{text}
========================================================================
SCENARIO A -- WITH startup probe (boot 18s tolerated)
========================================================================
  t= 20s  startup probe: first success -> liveness/readiness probes begin
  t= 20s  endpoints: Pod ADDED TO Service rotation (readiness success)
  t= 70s  endpoints: Pod REMOVED FROM Service rotation (readiness failing -- no restart!)
  t= 95s  endpoints: Pod ADDED TO Service rotation (readiness success)
  t=150s  liveness probe: failure 1/3 (mode=deadlocked)
  t=160s  endpoints: Pod REMOVED FROM Service rotation (readiness failing -- no restart!)
  t=160s  liveness probe: failure 2/3 (mode=deadlocked)
  t=170s  liveness probe: failure 3/3 (mode=deadlocked)
  t=170s  RESTART #1 (deadlock detected)
  t=185s  startup probe: first success -> liveness/readiness probes begin
  t=185s  endpoints: Pod ADDED TO Service rotation (readiness success)
------------------------------------------------------------------------
  restarts                    : 1
  seconds receiving traffic   : 171 / 240
  ready at end of simulation  : True

========================================================================
SCENARIO B -- WITHOUT startup probe (liveness doubles as boot check, boot 45s)
========================================================================
  t= 10s  liveness probe: failure 1/3 (mode=booting)
  t= 20s  liveness probe: failure 2/3 (mode=booting)
  t= 30s  liveness probe: failure 3/3 (mode=booting)
  t= 30s  RESTART #1 (killed while still booting)
  t= 40s  liveness probe: failure 1/3 (mode=booting)
  t= 50s  liveness probe: failure 2/3 (mode=booting)
  t= 60s  liveness probe: failure 3/3 (mode=booting)
  t= 60s  RESTART #2 (killed while still booting)
  t= 70s  liveness probe: failure 1/3 (mode=booting)
  t= 80s  liveness probe: failure 2/3 (mode=booting)
  t= 90s  liveness probe: failure 3/3 (mode=booting)
  t= 90s  restart budget exhausted -> CrashLoopBackOff; kubelet backs off,
          Pod never becomes ready
------------------------------------------------------------------------
  restarts                    : 3
  seconds receiving traffic   : 0 / 240
  ready at end of simulation  : False
\end{minted}

Three readings. In Scenario~A the 35-second dependency outage produced a
\emph{routing} change (endpoint removed at $t{=}70$~s, three readiness
periods after the fault began) and \emph{zero} restarts --- correct,
because restarting cannot fix the database; the readiness
\texttt{failureThreshold: 3} is what keeps a single flaky probe from
flapping the endpoint. The deadlock, a process-local fault, earned the
only restart, at $t{=}170$~s --- three liveness periods after it began.
In Scenario~B the \emph{identical} liveness configuration is fatal for a
different reason: a 45-second boot exceeds the $10\,\text{s} \times 3$
budget, so the kubelet kills the container at $t{=}30$~s, mid-boot,
forever. The configuration did not change between scenarios; the
\emph{meaning} of the liveness probe did. That is why probes are a
design artifact and not boilerplate.

\subsection{Listing 17.6 --- Canary Analyzer (\texttt{canary\_analysis.py})}
\label{sec:ch17-canary}
\label{lst:ch17-canary}

The analyzer consumes two metric series --- a long baseline and a
batched canary --- and makes the promote/rollback call with Wald's
Sequential Probability Ratio Test. For Bernoulli error events with
hypotheses $H_0: p = p_0$ (baseline error rate) and $H_1: p = p_1 = p_0 +
\delta$ (an intolerable regression), the log-likelihood ratio accumulates
per observation:
\[
\Lambda_n \;=\; \sum_{i=1}^{n} \Big[ x_i \ln\frac{p_1}{p_0} \;+\; (1 -
x_i)\,\ln\frac{1-p_1}{1-p_0} \Big], \qquad
\Lambda_n \ge \ln\tfrac{1-\beta}{\alpha} \Rightarrow \text{ROLLBACK},
\quad
\Lambda_n \le \ln\tfrac{\beta}{1-\alpha} \Rightarrow \text{PROMOTE},
\]
with $\alpha$ the false-promote budget and $\beta$ the false-rollback
budget (both 0.05 here, so both boundaries sit at $\pm\ln 19 \approx
\pm 2.944$). The SPRT's virtue for rollouts is that it is
\emph{sequential}: it decides as soon as the evidence allows --- a badly
regressed canary is caught in two or three batches instead of after the
full window --- with controlled error rates, unlike repeated peeking at a
fixed-horizon test. A hard latency gate (batch p95 $\le$ baseline p95
$\times 1.25$) runs alongside, because some regressions are pure latency.

\begin{minted}{python}
"""Canary analysis: error-rate delta + Wald SPRT promote/rollback decision.

Given a stable baseline metric series and a canary metric series (synthetic
fixtures generated with fixed seeds), this script computes, per analysis
batch:

* the canary/baseline error-rate delta,
* a Wald Sequential Probability Ratio Test (SPRT) on Bernoulli error
  events, which decides PROMOTE (canary no worse) or ROLLBACK (canary
  regressed by at least delta) as soon as the evidence crosses a boundary,
* a hard latency gate: batch p95 must stay within 25% of baseline p95.

Stdlib only, fully deterministic: every run prints identical output.
"""

from __future__ import annotations

import math
import random
from dataclasses import dataclass

# Configuration (the numbers an SRE would set in the rollout policy)
BATCH_SIZE = 200           # canary requests analyzed per gate evaluation
MAX_BATCHES = 10           # give up (hold) if still undecided after this many
DELTA = 0.010              # H1: canary error rate worse by >= 1 percentage pt
ALPHA = 0.05               # false-promote probability budget
BETA = 0.05                # false-rollback probability budget
LATENCY_P95_FACTOR = 1.25  # hard gate: canary p95 <= baseline p95 * 1.25


@dataclass(frozen=True)
class Event:
    latency_ms: float
    error: bool


def synth_series(n: int, mean_latency: float, jitter: float,
                 error_rate: float, seed: int) -> list[Event]:
    """Deterministic synthetic request series (lognormal-ish latencies)."""
    rng = random.Random(seed)
    series: list[Event] = []
    for _ in range(n):
        latency = mean_latency + 10.0 * rng.lognormvariate(0.0, jitter)
        series.append(Event(round(latency, 1), rng.random() < error_rate))
    return series


def percentile(sorted_values: list[float], pct: float) -> float:
    """Nearest-rank percentile; input must already be sorted ascending."""
    if not sorted_values:
        raise ValueError("percentile of empty series")
    rank = max(1, math.ceil(pct / 100.0 * len(sorted_values)))
    return sorted_values[rank - 1]


class BernoulliSPRT:
    """Wald's sequential test: H0 (p = p0) vs H1 (p = p1 = p0 + delta)."""

    def __init__(self, p0: float, p1: float, alpha: float, beta: float) -> None:
        if not (0.0 < p0 < p1 < 1.0):
            raise ValueError("require 0 < p0 < p1 < 1")
        self.p0, self.p1 = p0, p1
        self.upper = math.log((1.0 - beta) / alpha)   # rollback boundary
        self.lower = math.log(beta / (1.0 - alpha))   # promote boundary
        self.llr = 0.0
        self.errors = 0
        self.total = 0

    def observe(self, error: bool) -> None:
        self.total += 1
        if error:
            self.errors += 1
            self.llr += math.log(self.p1 / self.p0)
        else:
            self.llr += math.log((1.0 - self.p1) / (1.0 - self.p0))

    def verdict(self) -> str:
        if self.llr >= self.upper:
            return "ROLLBACK"
        if self.llr <= self.lower:
            return "PROMOTE"
        return "CONTINUE"


def analyze(label: str, baseline: list[Event], canary: list[Event]) -> str:
    base_errors = sum(1 for e in baseline if e.error)
    p0 = base_errors / len(baseline)
    base_lat = sorted(e.latency_ms for e in baseline)
    base_p95 = percentile(base_lat, 95)
    p95_limit = base_p95 * LATENCY_P95_FACTOR

    print("=" * 74)
    print(f"FIXTURE: {label}")
    print(f"baseline: {len(baseline)} reqs, {base_errors} errors "
          f"({p0:.2%}), p50={percentile(base_lat, 50):.1f} ms, "
          f"p95={base_p95:.1f} ms")
    print(f"H0: p_canary = {p0:.2%}   H1: p_canary >= {p0 + DELTA:.2%}   "
          f"alpha=beta={ALPHA}   latency gate: p95 <= {p95_limit:.1f} ms")
    print("-" * 74)
    print(f"{'batch':>5} {'reqs':>5} {'errors':>6} {'err-rate':>9} "
          f"{'delta':>8} {'p95 ms':>8} {'lat-gate':>8} {'LLR':>8}  decision")
    print("-" * 74)

    sprt = BernoulliSPRT(p0, p0 + DELTA, ALPHA, BETA)
    final = "HOLD (inconclusive after max batches)"
    for batch_no in range(1, MAX_BATCHES + 1):
        lo, hi = (batch_no - 1) * BATCH_SIZE, batch_no * BATCH_SIZE
        batch = canary[lo:hi]
        if not batch:
            break
        latency_breach = False
        for event in batch:
            sprt.observe(event.error)
        batch_lat = sorted(e.latency_ms for e in batch)
        batch_p95 = percentile(batch_lat, 95)
        if batch_p95 > p95_limit:
            latency_breach = True
        err_rate = sprt.errors / sprt.total
        delta = err_rate - p0
        verdict = sprt.verdict()
        print(f"{batch_no:>5} {sprt.total:>5} {sprt.errors:>6} "
              f"{err_rate:>9.2%} {delta:>+8.2%} {batch_p95:>8.1f} "
              f"{'BREACH' if latency_breach else 'ok':>8} "
              f"{sprt.llr:>8.3f}  {verdict}")
        if latency_breach:
            final = f"ROLLBACK (latency gate breached at batch {batch_no})"
            break
        if verdict == "ROLLBACK":
            final = (f"ROLLBACK (error regression detected at batch "
                     f"{batch_no}, {sprt.total} canary requests)")
            break
        if verdict == "PROMOTE":
            final = (f"PROMOTE (error parity proven at batch {batch_no}, "
                     f"{sprt.total} canary requests)")
            break
    print("-" * 74)
    print(f"FINAL DECISION: {final}")
    print()
    return final


def main() -> None:
    print("Canary analysis -- Northwind Robotics Orders API rollout")
    print(f"policy: batches of {BATCH_SIZE}, error delta >= {DELTA:.1%} "
          f"rolls back, p95 latency x{LATENCY_P95_FACTOR} rolls back")
    print()

    baseline = synth_series(4000, mean_latency=45.0, jitter=0.6,
                            error_rate=0.004, seed=2026)

    good = synth_series(MAX_BATCHES * BATCH_SIZE, mean_latency=47.0,
                        jitter=0.6, error_rate=0.004, seed=1701)
    analyze("canary v2.4.1 -- healthy (error parity, latency parity)",
            baseline, good)

    bad = synth_series(MAX_BATCHES * BATCH_SIZE, mean_latency=58.0,
                       jitter=0.6, error_rate=0.014, seed=1702)
    analyze("canary v2.4.2 -- regressed (error rate 3.5x baseline)",
            baseline, bad)


if __name__ == "__main__":
    main()
\end{minted}

Observed output:

\begin{minted}{text}
FIXTURE: canary v2.4.1 -- healthy (error parity, latency parity)
baseline: 4000 reqs, 19 errors (0.47%), p50=54.9 ms, p95=71.2 ms
H0: p_canary = 0.47%   H1: p_canary >= 1.47%   alpha=beta=0.05   latency gate: p95 <= 89.0 ms
--------------------------------------------------------------------------
batch  reqs errors  err-rate    delta   p95 ms lat-gate      LLR  decision
--------------------------------------------------------------------------
    1   200      1     0.50%   +0.03%     76.5       ok   -0.877  CONTINUE
    2   400      1     0.25%   -0.22%     70.7       ok   -2.896  CONTINUE
    3   600      2     0.33%   -0.14%     74.4       ok   -3.773  PROMOTE

FINAL DECISION: PROMOTE (error parity proven at batch 3, 600 canary requests)

FIXTURE: canary v2.4.2 -- regressed (error rate 3.5x baseline)
(baseline and policy header lines identical to the first fixture)
batch  reqs errors  err-rate    delta   p95 ms lat-gate      LLR  decision
    1   200      2     1.00%   +0.53%     88.0       ok    0.267  CONTINUE
    2   400      5     1.25%   +0.78%     88.7       ok    1.677  CONTINUE
    3   600     11     1.83%   +1.36%     85.9       ok    6.516  ROLLBACK

FINAL DECISION: ROLLBACK (error regression detected at batch 3, 600 canary requests)
\end{minted}

Both fixtures resolve in three batches --- 600 canary requests --- but in
opposite directions, and the trace shows why the sequential test is the
right tool: batch~1 of the regressed canary ($+0.53$ points of delta,
LLR $0.267$) is exactly the evidence a dashboard-watcher would shrug at
(``could be noise''), and batch~3 crosses the rollback boundary with the
error budget still honored. Note also what the analyzer deliberately
leaves to you: the latency gate never fired, because the regressed
fixture's mean-latency rise lands just under the gate
($88.0 \le 89.0$~ms) --- retune \texttt{LATENCY\_P95\_FACTOR} to
\texttt{1.20} in Exercise~7.4 and the same fixture falls to the latency
gate at the very first batch, two batches before the error SPRT speaks.

%% ----------------------------------------------------------------------------
\section{Common Pitfalls \& Anti-Patterns}
\label{sec:ch17-pitfalls}

\subsection*{\texttt{:latest} tags in anything that matters}
\texttt{latest} is not a version; it is a pointer that moves when the
publisher rebuilds. Deployments referencing it are unreproducible (two
clusters running ``the same'' manifest run different bytes), rollbacks
are impossible (there is no previous pointer to return to), and
Kubernetes' default \texttt{imagePullPolicy: IfNotPresent} means two
nodes can run \emph{different} images under one tag after a republish.
Pin tags to release versions; pin digests for anything you would have to
defend in an incident review.

\subsection*{Root containers and writable filesystems}
``It needs root'' almost always means ``we never checked.'' A FastAPI
service needs to bind a high port, read its own files, and write to one
scratch directory --- none of which is UID~0. Running as root with a
writable root filesystem converts every remote-code-execution bug into a
full container takeover with persistence (dropped binaries, modified
config). \texttt{USER} in the image, \texttt{runAsNonRoot} in the pod,
\texttt{readOnlyRootFilesystem: true}, \texttt{drop: ["ALL"]}
capabilities: four lines, most of a container-escape kill chain.

\subsection*{Liveness = readiness copy-paste}
Duplicating the readiness probe's path into the liveness probe is the
single most common probe bug in the wild, and it inverts the recovery
theory of Table~\ref{tab:ch17-probes}: a dependency outage (readiness
correctly fails) now also restarts every pod simultaneously, which
hammers the recovering dependency with reconnection storms and produces
a fleet-wide outage from a partial one. Rule of thumb: the liveness
handler should be \emph{cheaper} than the readiness handler and touch
\emph{fewer} things --- ideally just ``event loop turns.''

\subsection*{CPU limits causing throttle-induced latency}
Section~\ref{sec:ch17-k8s} derived the mechanism; the anti-pattern is the
cargo-culted \texttt{limits.cpu} copied from a tutorial. The signature:
p99 latency glued to multiples of ~100~ms, average CPU well below the
limit, \texttt{container\_cpu\_cfs\_throttled\_periods\_total} climbing.
The fix is not ``raise the limit'' blindly; it is ``decide whether the
limit's protection (noisy-neighbor containment) is worth its cost (tail
latency) for \emph{this} service'' --- latency-sensitive APIs usually
drop the CPU limit and rely on requests plus HPA.

\subsection*{HPA without a PDB (and PDBs without thought)}
An HPA will happily scale to \texttt{minReplicas} during a node drain,
and without a PDB the drain evicts \emph{all} of them --- autoscaling
governs desired count, drains govern eviction, and only the PDB sits
between them. The symmetric bug: \texttt{minAvailable: <replicas>} or
\texttt{maxUnavailable: 0} on a 1-replica deployment, which makes every
cluster upgrade hang until a human notices. PDBs are arithmetic, not
decoration: budget $=$ replicas you can lose \emph{and still serve the
SLO}, minus one for the rollout window.

\subsection*{Hardcoded secrets in manifests}
\texttt{DATABASE\_URL: postgres://user:password@\ldots} in a committed
YAML file is a credential in version control --- replicated to every
clone, every CI log that echoes the apply, every \texttt{kubectl get -o
yaml}. \texttt{secretKeyRef} is the floor; external-secrets operators or
CSI secret stores are the grown-up version. Also: \texttt{kubectl
create secret} values are base64, not encryption --- enable encryption at
rest and RBAC accordingly (Chapter~\ref{ch:api_security}).

\subsection*{TLS terminated once, plaintext to the pod}
``The LB does TLS'' plus plaintext HTTP from LB to pod means every
internal hop --- node network, mesh-less sidecar gap, packet mirror ---
carries bearer tokens in the clear, and a single misconfigured security
group exposes the service port directly. Defense in depth is cheap here:
re-encrypt (edge cert outside, internal CA inside) or run a mesh and let
mTLS cover east-west by default (Section~\ref{sec:ch17-edge}).

\subsection*{Migration deployed before compatible code}
The one rollout mistake rollback cannot fix: a destructive migration
(drop/rename) applied while old binaries still run, followed by a failed
rollout that restores old code against a schema it no longer
understands. Expand-and-contract exists precisely to make schema and
code independently deployable (Chapter~\ref{ch:versioning}): expand
first, deploy code that tolerates both shapes, contract only after every
old pod is gone. Pipeline-enforced, not remembered.

%% ----------------------------------------------------------------------------
\section{Hands-On Production Lab \& Real-World Case Study}
\label{sec:ch17-lab}

\subsection{Lab: Signals, Probes, and a Manifest Design Review}
Time: 90 minutes. Requirements: the venv from
Section~\ref{sec:ch17-objectives}; no Docker or cluster needed for parts
1--2.

\textbf{Part 1 --- Prove the drain (30 min).} Save Listings 17.3 and 17.4
side by side, then run \texttt{python shutdown\_demo.py}. Confirm all
three PASS lines, then break things deliberately: (a)~edit
\texttt{orders\_server.py} to add
\texttt{timeout\_graceful\_shutdown=1} to the \texttt{uvicorn.run} call
and re-run --- the export now fails mid-flight (act~3's timeout
amputating act~2's drain); (b)~restore it, and note that no edit can
make the drain work if the process never receives the signal --- on
Linux, wrap the entrypoint in a shell (\texttt{sh -c "uvicorn \ldots"})
and \texttt{SIGTERM} dies in the shell. Record both failure signatures.

\textbf{Part 2 --- Replay the probes (25 min).} Run
\texttt{python probe\_simulator.py}. Then edit the timeline: move the
\texttt{DEP\_DOWN} window to overlap the deadlock, and predict ---
before running --- the restart count and traffic seconds. Now set
\texttt{READINESS.failure\_threshold = 1} and observe the endpoint
flapping during the dependency outage (removed at the first failure,
oscillating with the probe period). Explain, in one sentence each, why
flapping hurts (connection churn, LB convergence cost) and why
threshold~3 was the right default.

\textbf{Part 3 --- Manifest design review (35 min).} Print
Listing~\ref{lst:ch17-deployment}. Walk it as a design-review meeting,
field by field, answering aloud: What happens to in-flight requests when
this pod is evicted? When the node dies? When the database is down for
a minute? What is the worst-case traffic loss during a rollout
(\texttt{maxUnavailable: 0} $\Rightarrow$ none below 4 ready) and what
does it cost (quota for one surge pod)? Why does CPU have no limit?
Where does the database password live, and who can read it? If any
answer is ``I don't know,'' the manifest has a bug you haven't found
yet --- keep digging until every field has a failure story.

\subsection{Case Study: The Rolling Deploy That Dropped In-Flight Requests}
\label{sec:ch17-case}

\textbf{Context.} Northwind Robotics' checkout-adjacent Orders API: 4
replicas behind a cloud load balancer, Kubernetes rolling updates,
\texttt{maxSurge: 1, maxUnavailable: 0} --- a configuration review had
declared it ``zero-downtime safe.'' Every Friday deploy nonetheless
produced a 30--60-second spike of HTTP~502s at exactly the rollout
minutes, roughly 0.1\% of requests --- small enough to lose in the
graphs, large enough to be a standing client complaint.

\textbf{Diagnosis.} The 502s correlated with pod terminations, and the
terminated pods' logs showed active requests severed mid-flight. The
mechanism was the race of Section~\ref{sec:ch17-process}: the kubelet
sent \texttt{SIGTERM} and \emph{concurrently} updated endpoints, but the
cloud load balancer's target list converged seconds later --- during
which it kept opening new connections to a process that had already
stopped accepting. Worse, the app's longest endpoint (the order export)
ran up to 40~s while \texttt{terminationGracePeriodSeconds} was the
default 30~s: even requests that \emph{were} in flight at SIGTERM were
SIGKILLed at the 30-second wall. Two gaps, one symptom. There was no
\texttt{preStop} hook at all --- the missing line the whole outage hung
on.

\textbf{Fix.} Three manifest lines and one server flag:
\texttt{preStop.sleep.seconds: 8} (the pod keeps serving while endpoint
removal fans out), \texttt{terminationGracePeriodSeconds: 60}
(slowest request 40~s $+$ preStop 8~s $+$ margin), uvicorn's graceful
timeout raised to match, and an SLO alert on 502 rate so the class of
bug could never again hide at 0.1\%. The next deploy's 502 count was
zero, and has been since.

\textbf{Lessons.} (1)~``Zero downtime'' is a property of the
\emph{termination path}, not the update strategy --- maxSurge arithmetic
protects capacity, not connections. (2)~Grace periods are budgets to be
engineered from measured request durations, not defaults to be
inherited. (3)~The correct test is exactly Listing~\ref{lst:ch17-shutdown}:
signal the server mid-request and watch from the client side; had that
test existed, the config review's ``safe'' stamp would have died in CI
instead of in production.

%% ----------------------------------------------------------------------------
\section{Self-Assessment Exercises \& Deep-Dive Questions}
\label{sec:ch17-exercises}

\begin{enumerate}[label=\textbf{Exercise 7.\arabic*}, leftmargin=2.9cm]
  \item \textbf{Layer forensics.} Rewrite Listing 17.1's Dockerfile so
        the source \texttt{COPY} precedes the dependency install.
        Explain exactly which layers rebuild on a one-line source edit
        now, and quantify the cost for a 200~MB dependency layer.
  \item \textbf{Break the drain.} In \texttt{orders\_server.py}, set
        \texttt{timeout\_graceful\_shutdown=1} and capture the demo's new
        failure signature. Which of the three verdict lines flips, and
        which act of the four-act shutdown did you just amputate?
  \item \textbf{Probe arithmetic.} A service boots in 75 seconds. Using
        \texttt{periodSeconds: 5}, what is the minimum
        \texttt{failureThreshold} for a startup probe that tolerates this
        with 25\% headroom? What goes wrong if you instead raise the
        liveness \texttt{failureThreshold} to accommodate it?
  \item \textbf{Retune the gate.} Run \texttt{canary\_analysis.py} with
        \texttt{LATENCY\_P95\_FACTOR = 1.20} and confirm the regressed
        fixture is rolled back at batch~1 by the latency gate. Then set
        \texttt{DELTA = 0.02}: what happens to the regressed fixture's
        detection latency, and what does that teach about choosing
        $\delta$ relative to your error budget?
  \item \textbf{CrashLoop replay.} In \texttt{probe\_simulator.py},
        give Scenario~B an 18-second boot (same as A). Does it still
        crash-loop? Now 35 seconds. Derive the exact boot-time boundary
        between the two outcomes from the probe parameters alone.
  \item \textbf{Endpoint lag budget.} Your cloud LB converges target
        removals in up to 12~s and your p99.9 request duration is 9~s.
        Choose \texttt{preStop} sleep, uvicorn graceful timeout, and
        \texttt{terminationGracePeriodSeconds}, and defend each number.
  \item \textbf{Design review script.} Take Listing
        \ref{lst:ch17-deployment} and remove the PDB. Write the exact
        sequence of events (drain, evict, HPA response) that takes the
        API fully offline during a routine node upgrade.
  \item \textbf{Canary slice sizing.} Baseline error rate 0.4\%,
        $\delta = 1$ point, $\alpha = \beta = 0.05$. Using the expected
        per-request LLR drift, estimate the expected number of canary
        requests to rollback when the true canary rate is 1.4\%, and to
        promote when it equals baseline. Verify with the simulator.
\end{enumerate}

\textbf{Deep-dive questions for further study:} How does
\texttt{containerd}'s image pull interact with digest pinning across a
node pool with heterogeneous pull times during a surge? What changes in
every section of this chapter when the mesh moves from sidecars to
ambient mode? How would you express Listing~\ref{lst:ch17-canary}'s gate
in PromQL-based analysis instead of batch files? What does the
\texttt{PodLifecycleSleepAction} feature gate's graduation history tell
you about depending on alpha features in lifecycle hooks?

%% ----------------------------------------------------------------------------
\section{Interview Q\&A Vault}
\label{sec:ch17-interview}

\subsection*{Junior (Q1--Q10)}
\begin{enumerate}[label=\textbf{Q\arabic*.}, leftmargin=1.6cm]
  \item \textbf{What is a container image, in one paragraph?}
        A content-addressed stack of filesystem layers plus a manifest:
        each layer is a diff addressed by its hash, the image is
        addressed by the manifest's digest, and a running container adds
        one thin writable layer on top. The properties that matter are
        immutability (bytes never change after build) and traceability
        (the digest names exactly one build forever).
  \item \textbf{What problem does a multi-stage Dockerfile solve?}
        The runtime image inherits everything the build needed ---
        compilers, package managers, caches --- unless you build in one
        stage and copy only artifacts into a fresh final stage. The win
        is triple: smaller images (faster pulls, faster cold starts),
        smaller attack surface (fewer binaries for an attacker to reuse),
        and cleaner caching (dependency layers invalidate only when
        manifests change).
  \item \textbf{Why should containers not run as root?}
        Container root is a UID on the host separated only by namespace
        configuration; a container-escape or a misconfigured mount turns
        application-root into something dangerously close to host-root.
        Running as a fixed unprivileged UID forces an attacker to win
        twice: exploit the app, \emph{then} escalate. It is also a
        deploy-time assertion: \texttt{runAsNonRoot: true} makes
        Kubernetes refuse an image that regresses.
  \item \textbf{What does \texttt{.dockerignore} protect?}
        Two things: build correctness (junk churn like \texttt{.git} or
        local logs would otherwise bust the cache on every build) and
        secrets (a stray \texttt{.env} copied by \texttt{COPY . .}
        persists in the layer history even if a later layer deletes it).
  \item \textbf{What is the difference between a tag and a digest?}
        A tag (\texttt{:3.12-slim}) is a mutable registry pointer that
        can be retargeted at any time; a digest
        (\texttt{@sha256:\ldots}) is the content hash of one manifest and
        can never name anything else. Tags are for humans; digests are
        for deployments you may have to defend in an incident review.
  \item \textbf{What happens, in order, when Kubernetes terminates a
        Pod?}
        The Pod is marked terminating and removed from endpoints
        (asynchronously); the \texttt{preStop} hook runs; then
        \texttt{SIGTERM} to PID~1; the grace period
        (\texttt{terminationGracePeriodSeconds}, default 30~s) counts
        down; then \texttt{SIGKILL}. Note the trap: endpoint removal is
        concurrent, not first --- traffic can still arrive after
        \texttt{SIGTERM}.
  \item \textbf{What is a liveness probe?}
        A periodic check (HTTP, TCP, or exec) whose consecutive failures
        --- \texttt{failureThreshold} of them --- make the kubelet
        \emph{restart the container in place}. It answers one question:
        is this process unrecoverably stuck? It should therefore check
        only process-local health and never a dependency.
  \item \textbf{What is an API gateway?}
        A reverse proxy at the platform edge that owns cross-cutting
        north-south policy --- routing, JWT validation offload, rate
        limiting, transformation, canary splitting --- so those concerns
        are configured once instead of reimplemented per service.
  \item \textbf{What is blue-green deployment?}
        Two complete environments; traffic switches atomically from one
        to the other (selector, LB target, or DNS), and rollback is the
        switch flipped back. Instant rollback and a warm target, at the
        price of double infrastructure and an idle environment that rots
        unless exercised.
  \item \textbf{What is a PodDisruptionBudget?}
        A policy object stating how many pods of a service may be
        \emph{voluntarily} unavailable (\texttt{minAvailable} or
        \texttt{maxUnavailable}); drains and cluster upgrades respect it
        by blocking, while crashes (involuntary) ignore it. Without one,
        a routine node drain can evict your entire fleet at once.
\end{enumerate}

\subsection*{Mid (Q11--Q20)}
\begin{enumerate}[label=\textbf{Q\arabic*.}, leftmargin=1.6cm, start=11]
  \item \textbf{Liveness vs.\ readiness vs.\ startup probes --- failure
        behavior of each?}
        Startup: only probe evaluated during boot; exhausting
        \texttt{periodSeconds} $\times$ \texttt{failureThreshold}
        restarts the container \emph{before} it ever served --- it exists
        to give slow starters a budget liveness never sees. Readiness:
        failure removes the Pod from Service endpoints --- a routing
        consequence, no restart --- and the first success re-adds it; it
        answers ``should I get traffic \emph{now}?'' Liveness:
        \texttt{failureThreshold} consecutive failures restart the
        container in place; it answers ``am I dead in a way only a
        restart can fix?'' The simulator in
        Section~\ref{sec:ch17-probesim} demonstrates all three
        consequences over one fault timeline.
  \item \textbf{Why is copying the readiness probe into liveness
        dangerous?}
        Readiness legitimately fails for external reasons (dependency
        down); making liveness fail for the same reason converts a
        partial outage into a fleet-wide restart storm --- every pod
        restarts simultaneously and then reconnection-storms the
        dependency that was trying to recover. Readiness says
        ``wait''; liveness says ``kill me''; merging them deletes the
        distinction that recovery depends on.
  \item \textbf{Why can CPU limits \emph{hurt} latency?}
        CPU limits are enforced by CFS quota throttling in (default)
        100~ms periods: consume the quota in the first few milliseconds
        of a burst and every thread in the container freezes until the
        next period. The signature is p99 latency quantized to ~100~ms
        multiples while average CPU sits far below the limit. Memory
        limits OOM-kill instead --- violent but honest. Mitigation:
        requests always (scheduling and CFS shares), CPU limits omitted
        or generous on latency-sensitive services, and an alert on the
        throttled-periods metric.
  \item \textbf{Explain \texttt{maxSurge} and \texttt{maxUnavailable}.}
        Rolling-update trade-off knobs. \texttt{maxSurge}: extra pods
        above the replica count allowed during rollout (needs quota).
        \texttt{maxUnavailable}: replicas allowed to be down. With
        $R{=}4$: \texttt{surge 1 / unavailable 0} never serves below 4
        and costs one extra pod; \texttt{surge 0 / unavailable 1} costs
        nothing but serves a window at 3. \texttt{0/0} cannot roll at
        all --- a common copy-paste wedge.
  \item \textbf{What is endpoint-removal lag during scale-down?}
        ``Pod terminating'' and ``endpoints object updated'' are
        concurrent, and downstream convergence --- kube-proxy iptables
        programming, DNS TTLs, cloud LB target lists, gateway discovery
        --- takes seconds more. So a pod can stop accepting (post-SIGTERM)
        while the LB still opens new connections to it: instant 502s.
        Mitigation: \texttt{preStop} sleep so the process keeps serving
        through the lag, plus client retries for the residue
        (Chapter~\ref{ch:microservices}).
  \item \textbf{How does an HPA compute replicas, and what is the
        stabilization window?}
        Every sync period:
        $\lceil \text{current} \times \text{currentMetric} /
        \text{targetMetric} \rceil$, where resource utilization is
        measured against \emph{requests} (inflated requests silently
        disable autoscaling). The scale-down stabilization window
        (default 300~s) takes the \emph{highest} recommendation in the
        window, because every scale-down terminates connections; scale-up
        is typically immediate. Downscale policies cap pods removed per
        period.
  \item \textbf{Design a canary release with automated rollback --- the
        sketch?}
        Small traffic steps (1/5/25/50/100\% or request-budget batches),
        stable splitting (hash on API key/tenant so a client doesn't
        oscillate between versions), a pre-registered gate per step ---
        error-rate SPRT with budgets $\alpha$/$\beta$ plus a hard
        latency-quantile guardrail, as in
        Listing~\ref{lst:ch17-canary} --- and rollback defined as the
        same weight change in reverse. Prerequisites: expand-and-contract
        schema so old code still reads new data, and per-version metric
        labels so the gate can see the canary at all.
  \item \textbf{Gateway vs.\ ingress vs.\ service mesh?}
        Ingress is the Kubernetes-standard vocabulary for north-south
        HTTP routing plus TLS termination --- a gateway subset. An API
        gateway is the product tier at the platform edge: auth offload,
        rate limiting, transformation, analytics --- the policy
        enforcement point for north-south traffic. A mesh puts proxies on
        every workload to govern \emph{east-west}: mTLS, retries,
        telemetry between services. Production commonly runs all three in
        sequence; the failure mode is enforcing one concern in two tiers
        with different budgets.
  \item \textbf{Where should TLS terminate for a public API?}
        At the provider edge with automated public certificates (ACME),
        then \emph{re-encrypt} to backends with an internal CA --- or let
        a mesh's mTLS cover it. Edge-only termination leaves the last hop
        plaintext (bearer tokens on the wire inside your VPC);
        passthrough gives end-to-end encryption but blinds the WAF and
        L7 routing entirely. The decision is: which tiers need to
        \emph{see} the traffic versus merely \emph{carry} it.
  \item \textbf{Why does a container need to care about PID~1?}
        The kernel delivers signals to PID~1 only if a handler is
        installed, and orphans re-parent to PID~1 for reaping. A shell
        wrapper as PID~1 swallows \texttt{SIGTERM} (the shell has no
        handler); Kubernetes then waits the grace period and SIGKILLs,
        severing in-flight work. Fix: exec-form \texttt{ENTRYPOINT} so
        the server \emph{is} PID~1, and a server (uvicorn) that installs
        real handlers.
\end{enumerate}

\subsection*{Senior (Q21--Q28)}
\begin{enumerate}[label=\textbf{Q\arabic*.}, leftmargin=1.6cm, start=21]
  \item \textbf{Walk me through every system that must cooperate for a
        request to survive a rolling deploy.}
        Endpoint controller removes the pod (async); preStop sleeps
        through LB convergence; SIGTERM; uvicorn stops accepting, drains
        within its graceful timeout; \texttt{terminationGracePeriodSeconds}
        exceeds preStop $+$ slowest request; LB draining stops new
        connections; clients retry the residue with backoff and
        idempotency keys (Chapter~\ref{ch:idempotency}); the new pod's
        startup probe tolerates cold boot; readiness gates its traffic.
        Any link missing and some fraction of requests dies.
  \item \textbf{Why is the SPRT better than ``watch the error rate for
        ten minutes''?}
        A fixed-horizon eyeball has uncontrolled error rates and no
        early stopping: a catastrophically bad canary burns ten minutes
        of damage, a healthy one waits ten minutes for nothing. The SPRT
        pre-registers hypotheses ($p_0$, $p_0{+}\delta$) and error
        budgets ($\alpha$, $\beta$), then decides the moment cumulative
        evidence crosses a boundary --- typically a few hundred requests
        either way --- with the false-promote and false-rollback rates
        \emph{guaranteed} by construction. It also forces the honest
        conversation upfront: what regression size $\delta$ do we
        actually care about?
  \item \textbf{Your p99 spiked after you ``hardened'' resources. No
        deploy. Diagnose.}
        The fingerprint: p99 quantized near 100~ms multiples, mean CPU
        far below limit, throttle counters climbing --- you added a CPU
        limit and CFS throttling is freezing all threads mid-request for
        the remainder of each quota period. Verify with
        \texttt{container\_cpu\_cfs\_throttled\_periods\_total}. Fix:
        drop the CPU limit (keep requests) or raise it well above p99
        burst usage; consider fewer worker processes so per-period quota
        lasts longer. Memory ``hardening'' would have shown OOMKills
        instead --- different graph, same meeting.
  \item \textbf{How do you size a canary's traffic steps?}
        By \emph{requests}, not percentages: the gate needs enough canary
        volume to detect the regression you care about. At baseline
        $p_0$ and intolerance $\delta$, expected requests-to-decide scale
        roughly with $1/\delta^2$; at $p_0{=}0.4\%$ and
        $\delta{=}1$~pt that is a few hundred requests per step
        (Listing~\ref{lst:ch17-canary} resolves in 600), while a 1\%
        slice of a low-traffic API might take hours --- so either accept
        coarser $\delta$, lengthen windows, or synthesize load. Step
        count trades blast radius against total rollout time.
  \item \textbf{Feature flags or canary deploys?}
        Different axes. Deployments change binaries; flags change code
        paths within one binary. Flags give per-tenant targeting and
        instant kill-switches but accrue combinatorial debt and live in
        the release path forever; canaries bound blast radius with zero
        code changes but run two contract versions simultaneously.
        Mature orgs compose them: canary the binary rollout, flag the
        risky feature \emph{within} the new binary, and never let a flag
        outlive its rollout.
  \item \textbf{The migration must ship with the code. Order them.}
        Expand-and-contract, never destructive-first: (1)~expand
        migration (add the new column/table; nothing renamed or dropped);
        (2)~deploy code that writes both shapes and reads tolerantly;
        (3)~backfill if needed; (4)~after \emph{every} old binary is
        gone, contract migration (drop the old shape). Violating the
        order makes rollback impossible: restored old code meets a schema
        it cannot read. This is Chapter~\ref{ch:versioning}'s discipline
        enforced by the deployment pipeline, not by memory.
  \item \textbf{Why does the Service selector in
        Listing~\ref{lst:ch17-service} omit the \texttt{track} label, and
        what does that buy and cost?}
        Omitting \texttt{track} makes one Service span stable and canary
        pods, splitting traffic by ready-pod ratio --- free canary
        weights with zero infrastructure. Cost: weights are quantized to
        pod counts (1/5, 2/6, \ldots), there is no session stability
        (one client can oscillate between versions), and autoscaling the
        stable track silently moves the split. Gateway weighted clusters
        (Envoy) fix all three at the price of a control plane.
  \item \textbf{What belongs in a readiness check --- and what must
        never be in a liveness check?}
        Readiness may shallowly check \emph{critical} dependencies
        (can I reach the database with a sub-100-ms probe?) because its
        consequence is only traffic removal. Liveness must check nothing
        external: its consequence is restart, and restarting cannot fix
        someone else's outage --- it only multiplies it. Deepest-check
        discipline: liveness $<$ readiness $<$ nothing at request time;
        a dependency-deep liveness probe is a cascading-failure machine
        (Chapter~\ref{ch:microservices}).
\end{enumerate}

\subsection*{Staff (Q29--Q36)}
\begin{enumerate}[label=\textbf{Q\arabic*.}, leftmargin=1.6cm, start=29]
  \item \textbf{Design the deployment safety system for a platform with
        200 API services.}
        Paved road, not gate: one blessed base image (distroless,
        digest-pinned, rebuilt on CVE), one manifest template with the
        probe/preStop/PDB/resource decisions pre-made and parameterized,
        progressive delivery as the \emph{only} rollout mechanism with
        platform-wide gates (error SPRT, latency guardrails, SLO burn
        rate), schema-change pipelines that physically cannot run
        contract before code, and automatic rollback as default with
        humans paged on gate \emph{failure to decide}, not on every
        rollout. Measure: change-failure rate, time-to-rollback, and
        percentage of services off-template.
  \item \textbf{When is a service mesh \emph{not} worth it?}
        When the east-west concerns it buys (mTLS, L7 retry policy,
        traffic telemetry, fine-grained splits) are already solved or not
        needed: few services, one team, flat compliance requirements, or
        a gateway already owning the splits you need. Costs are real: a
        proxy on every hop's latency budget, a control plane to operate,
        and a debugging story that now includes someone else's iptables.
        The 5-service monorepo does not need Istio; the 200-service
        regulated platform probably cannot live without it.
  \item \textbf{How do you make rollbacks boring?}
        Every property that makes rollbacks rare makes them safe:
        immutable digest-pinned artifacts (rollback is a pointer move,
        never a rebuild); expand-and-contract schemas (old code always
        reads current data); forward-compatible APIs
        (Chapter~\ref{ch:versioning}) so consumers tolerate either
        version; automated rollback triggers rehearsed on every deploy
        rather than saved for emergencies; and fast rollbacks measured
        and SLO'd like any other operation. The SRE framing: release
        engineering is risk management, and rollback is the cheapest
        risk control you own~\cite{beyer2016sre}.
  \item \textbf{A vendor gateway, NGINX, or Envoy-as-gateway for a new
        platform --- how do you decide?}
        Score on the matrix of Table~\ref{tab:ch17-gateway-matrix} plus
        two columns the table can't hold: operational ownership (who
        pages for it at 3~a.m.\ --- a vendor SLA or your team) and
        extensibility trajectory (plugin ecosystems and Wasm beat
        config-file reloads the moment policy gets bespoke). Default
        heuristic: if you will run a mesh within two years, standardize
        on Envoy-flavored everything now; if you need a developer portal
        and monetization, that's gateway-product territory (Kong class);
        if the traffic profile is simple and the team's depth is NGINX,
        boring wins.
  \item \textbf{How do you set error budgets $\alpha$, $\beta$, and
        $\delta$ for an automated canary gate?}
        From the SLO, backwards: $\delta$ is the regression size that
        would burn meaningful error budget if it reached 100\% for one
        detection window; $\alpha$ (false promote) is priced against
        rollback cost --- cheap rollback buys aggressive $\alpha$;
        $\beta$ (false rollback) is priced against developer velocity ---
        flaky gates train teams to bypass gates, which is how gates die.
        Then verify with the simulator: replay historical incident metric
        series through the gate and confirm it catches every real
        regression \emph{and} passes the last twenty healthy deploys.
        A gate that fails the replay is a gate you have not designed.
  \item \textbf{What is the blast-radius math behind progressive
        delivery?}
        Expected harm of a bad deploy $\approx$ (probability of
        regression) $\times$ (affected fraction) $\times$ (time to
        detect $+$ time to mitigate). Progressive delivery attacks all
        three factors: canary slices shrink the fraction, sequential
        gates shrink detection time from SLO-burn to seconds-of-evidence,
        and weight-based rollback shrinks mitigation to a pointer change.
        The complementary lever is shrink-the-unit: smaller, more
        frequent releases reduce the probability term, which is why
        release cadence and deployment safety are the same
        investment~\cite{nygard2018releaseit}.
  \item \textbf{How do you handle stateful connections (WebSockets,
        gRPC streams) under everything in this chapter?}
        Long-lived connections break every assumption above: draining is
        hours not seconds, L7 load balancing pins streams to pods, and
        HPA-on-CPU misses connection-count imbalance. The playbook:
        advertise \emph{max connection age} and make clients honor it
        (gRPC \texttt{MAX\_CONNECTION\_AGE}, WS renegotiation), so drains
        complete in bounded time; rebalance deliberately on scale events;
        size \texttt{terminationGracePeriodSeconds} to connection
        lifetime, not request latency; and consider a dedicated
        connection-tier deployment with its own, gentler rollout policy
        (Chapter~\ref{ch:async_apis}, Chapter~\ref{ch:grpc}).
  \item \textbf{Regulated industry: what changes in this chapter?}
        Everything gets evidence. Images: signed (cosign) and admitted
        by policy, digests everywhere, SBOMs per release. Probes and
        rollout params: change-controlled, reviewed like code. TLS:
        passthrough or re-encrypt with an internal CA --- plaintext
        last-hops stop being a judgment call. Secrets: external vault
        with rotation and audit, never cluster Secrets alone. Canary
        gates: thresholds become compliance artifacts, and the analyzer's
        decision log is a record you produce in audits. The engineering
        is the same; the difference is that ``why'' must be written down
        \emph{before} the deploy, not reconstructed after the incident.
  \item \textbf{What would you instrument to know, a year from now, that
        this chapter's discipline is working?}
        Leading indicators: percentage of deploys that are progressive
        with gates (target 100\%), rollback mean-time and rollback
        \emph{frequency} (both should be boring), gate false-positive
        rate (drives bypass behavior), probe-related incidents (flapping
        endpoints, boot-deadline crash loops --- target zero),
        CFS-throttling alerts on latency-critical services, deploy-time
        5xx rate (the case study's metric), and drift of running digests
        from declared releases. The meta-metric: whether the on-call
        can explain any of these numbers without opening a dashboard
        first.
\end{enumerate}

%% ----------------------------------------------------------------------------
\section{External Bibliography \& Further Reading}
\label{sec:ch17-bibliography}

\begin{itemize}
  \item \textbf{Kubernetes documentation --- Pod lifecycle and probes:}
        \emph{Container Probes} and \emph{Pod Lifecycle}
        (\url{https://kubernetes.io/docs/concepts/workloads/pods/pod-lifecycle/}).
        The authoritative semantics for startup/readiness/liveness
        behavior, termination ordering, and
        \texttt{terminationGracePeriodSeconds} that
        Section~\ref{sec:ch17-k8s} compresses.
  \item \textbf{Kubernetes documentation --- Horizontal Pod Autoscaler:}
        \url{https://kubernetes.io/docs/tasks/run-application/horizontal-pod-autoscale/}.
        The scaling algorithm, metric types, and the \texttt{behavior}
        stabilization windows and policies used in
        Listing~\ref{lst:ch17-hpa}.
  \item \textbf{Kubernetes documentation --- Disruptions:}
        \url{https://kubernetes.io/docs/concepts/workloads/pods/disruptions/}.
        Voluntary vs.\ involuntary disruption and PDB arithmetic ---
        required reading before your first node-pool upgrade.
  \item \textbf{Docker --- best practices for writing Dockerfiles:}
        \url{https://docs.docker.com/develop/develop-images/dockerfile_best-practices/}.
        Layer caching, instruction ordering, and multi-stage builds;
        pair with the \emph{Overview of best practices for building
        images} page for the attack-surface framing of
        Section~\ref{sec:ch17-image}.
  \item \textbf{Envoy proxy documentation:}
        \url{https://www.envoyproxy.io/docs/envoy/latest/}. Weighted
        clusters, the JWT authentication filter, local/global rate
        limiting, and the xDS APIs that make Envoy the substrate of
        modern gateways and meshes.
  \item \textbf{Kong Gateway documentation:}
        \url{https://docs.konghq.com/gateway/}. Plugin reference (JWT,
        rate-limiting, request-transformer, canary) and the declarative
        \texttt{deck} configuration model.
  \item \textbf{Google SRE book, chapter 8 --- \emph{Release
        Engineering}:} Betsy Beyer et al.\ (\cite{beyer2016sre}, also at
        \url{https://sre.google/sre-book/release-engineering/}). The
        canonical argument that releases are a risk-management
        discipline, not an afterthought.
  \item \textbf{Bilgin Ibryam \& Roland Hu\ss, \emph{Kubernetes
        Patterns} (O'Reilly, 2nd ed.):} the pattern vocabulary behind
        this chapter's manifests --- \emph{Predictable Demands}
        (requests/limits), \emph{Self Awareness}, health probes, and the
        lifecycle patterns that justify preStop and termination budgets.
  \item \textbf{The Twelve-Factor App}~\cite{twelvefactor}: factors III
        (config in the environment), IV (backing services as attached
        resources), and IX (disposability: fast startup, graceful
        shutdown) are this chapter's process discipline in manifesto
        form.
  \item \textbf{OWASP Cheat Sheet Series --- Docker Security and
        Kubernetes Security:}
        \url{https://cheatsheetseries.owasp.org/cheatsheets/Docker_Security_Cheat_Sheet.html}
        and the Kubernetes Security cheat sheet. The checklist form of
        Section~\ref{sec:ch17-pitfalls}'s container and manifest
        anti-patterns.
  \item \textbf{Michael Nygard, \emph{Release It!} (2nd
        ed.)}~\cite{nygard2018releaseit}: stability patterns and
        antipatterns whose production war stories motivate the
        connection-draining and rollout-safety sections.
  \item \textbf{Mark Richards \& Neal Ford's deployment chapters in
        \emph{Fundamentals of Software Architecture}} and Sam Newman's
        deployment discussion in \emph{Building
        Microservices}~\cite{newman2021microservices}: progressive
        delivery placed in its architectural context.
  \item \textbf{Wald, \emph{Sequential Analysis} (1947)} --- for the
        mathematically inclined, the original SPRT construction behind
        Listing~\ref{lst:ch17-canary}; a gentler modern treatment is any
        mathematical-statistics text's chapter on sequential testing.
\end{itemize}

%% ----------------------------------------------------------------------------
\section{Capstone Projects}
\label{sec:ch17-capstones}

\subsection*{Capstone 17.1 --- Harden the Image \hfill [junior]}
\textbf{Objective:} Take a naive single-stage Dockerfile for the Orders
API and rebuild it to this chapter's standard.

\textbf{Inputs/Datasets:} Listing 17.1; a ``naive'' Dockerfile you write
first (full base image, root, \texttt{COPY . .}, shell entrypoint); no
datasets.

\textbf{Steps:} (1)~Write the naive version and record its image size
and layer count. (2)~Apply, one commit at a time: multi-stage split,
slim or distroless runtime, non-root \texttt{USER} with fixed UID,
\texttt{.dockerignore}, cache-ordered \texttt{COPY}, \texttt{HEALTHCHECK},
exec-form \texttt{ENTRYPOINT}. (3)~Record size and (if a scanner is
available) finding count after each step. (4)~Produce a table of the
deltas with one sentence of justification per step.

\textbf{Acceptance Criteria:}
\begin{itemize}
  \item[$\square$] Final image is $\leq$ 20\% of the naive image's size
        (or you document precisely why not).
  \item[$\square$] \texttt{docker inspect} shows a non-root user and an
        exec-form entrypoint.
  \item[$\square$] A one-line source edit rebuilds only the source layer
        (build log evidence).
  \item[$\square$] No \texttt{.env}, \texttt{.git}, or test artifacts in
        any layer (\texttt{docker history} / \texttt{dive} evidence).
\end{itemize}

\textbf{Stretch Goals:} Produce both slim and distroless variants and
document the one operational capability (preStop \texttt{exec} sleep,
in-container debugging) you traded away with distroless.

\subsection*{Capstone 17.2 --- Prove the Drain Under Load \hfill [mid]}
\textbf{Objective:} Extend the graceful-shutdown demo from one in-flight
request to a realistic mix, and make the demo CI-grade.

\textbf{Inputs/Datasets:} Listings 17.3--17.4; synthetic traffic mix:
70\% fast \texttt{POST /orders}, 25\% \texttt{/healthz}, 5\% slow
\texttt{/orders/export}.

\textbf{Steps:} (1)~Rewrite the driver to run $N$ concurrent workers
generating the mix at a fixed request rate while the signal arrives at
a random-but-seeded moment. (2)~Count and classify outcomes: completed,
refused-after-signal (expected), and \emph{severed} (the failure class).
(3)~Sweep \texttt{timeout\_graceful\_shutdown} $\in \{0, 1, 10\}$~s and
tabulate severed counts. (4)~Wrap the whole experiment as a pytest that
asserts zero severed requests at the production timeout.

\textbf{Acceptance Criteria:}
\begin{itemize}
  \item[$\square$] Three verdict classes reported per run; totals
        reconcile exactly with requests attempted.
  \item[$\square$] Zero severed requests at the production timeout;
        nonzero at \texttt{timeout=0}; you can explain both numbers.
  \item[$\square$] Seeded and offline: two runs produce identical
        classification counts.
  \item[$\square$] pytest integration exits nonzero on any severed
        request.
\end{itemize}

\textbf{Stretch Goals:} Add a preStop-equivalent delay to the driver
(keep serving $k$ seconds post-signal) and show the refused-after-signal
count drop to zero.

\subsection*{Capstone 17.3 --- Manifest Review Board \hfill [mid]}
\textbf{Objective:} Run Listing 17.2's manifest set through a formal
design review and fix what you find.

\textbf{Inputs/Datasets:} Listing~\ref{lst:ch17-deployment} and its
companions, saved as YAML files; \texttt{pyyaml} for the linting script;
a seeded list of injected faults (you create: remove the PDB, set
\texttt{minAvailable: 4}, add \texttt{limits.cpu: 100m}, delete the
preStop hook, hardcode \texttt{DATABASE\_URL}, set
\texttt{terminationGracePeriodSeconds: 5}, copy readiness into liveness).

\textbf{Steps:} (1)~Write \texttt{manifest\_lint.py}: parse each YAML
document and assert this chapter's rules (probes distinct paths; memory
request $=$ limit; no CPU limit or a justified one; PDB present and
arithmetic consistent with replicas; no literal secret-looking strings;
preStop present; grace period $>$ slowest endpoint budget).
(2)~Verify the linter passes the clean set and catches every injected
fault, naming the rule violated. (3)~Write the review minutes: for each
caught fault, the failure story it prevents.

\textbf{Acceptance Criteria:}
\begin{itemize}
  \item[$\square$] Linter catches all seven injected faults with
        rule-specific messages.
  \item[$\square$] Linter passes the clean manifest set.
  \item[$\square$] Every lint rule cites its chapter section.
  \item[$\square$] Review minutes contain a failure story per fault, not
        restatements of the rule.
\end{itemize}

\textbf{Stretch Goals:} Add a second reviewer mode that diffs two
manifest revisions and flags only \emph{safety-relevant} changes
(probes, resources, lifecycle, grace periods).

\subsection*{Capstone 17.4 --- Gate Tuner \hfill [senior]}
\textbf{Objective:} Turn Listing~\ref{lst:ch17-canary} from a demo into
a defensible rollout policy.

\textbf{Inputs/Datasets:} \texttt{canary\_analysis.py}; a seeded fixture
generator you extend with a third case: a canary whose regression is
exactly $0.5\delta$ (below the intolerance line --- the gate
\emph{should} promote it, eventually).

\textbf{Steps:} (1)~Sweep $\delta \in \{0.5, 1.0, 2.0\}$ points and
$\alpha = \beta \in \{0.01, 0.05, 0.10\}$ over the three fixtures;
record requests-to-decision and verdict for each cell. (2)~Add a
latency SPRT (Gaussian approximation on per-batch means) alongside the
hard p95 gate and compare detection times. (3)~Derive, on paper, the
expected LLR drift per request as a function of the true canary rate,
and verify the simulator's detection latencies match within a batch.
(4)~Write the one-page rollout policy the numbers justify.

\textbf{Acceptance Criteria:}
\begin{itemize}
  \item[$\square$] Sweep table shows the $\delta$/$\alpha$/$\beta$
        trade-off: tighter intolerance costs more requests per decision.
  \item[$\square$] The $0.5\delta$ canary is promoted (never rolled
        back) in every cell.
  \item[$\square$] Measured detection latency matches the drift
        derivation within one batch size.
  \item[$\square$] The policy page names $\delta$, $\alpha$, $\beta$,
        batch size, and the rollback action --- and cites the evidence
        for each.
\end{itemize}

\textbf{Stretch Goals:} Add a seasonal-baseline mode: compare the canary
against a time-shifted baseline (same hour last week) and quantify what
seasonality does to the SPRT's assumptions.

\subsection*{Capstone 17.5 --- Zero-Downtime Reference Architecture
\hfill [staff]}
\textbf{Objective:} Design and document the complete deployment
architecture for Northwind Robotics' public API platform --- the chapter,
applied end to end.

\textbf{Inputs/Datasets:} All chapter listings; the platform brief: 12
services (3 latency-critical), mixed REST and gRPC, one WebSocket
telemetry service, PCI-adjacent payments path, two regions.

\textbf{Steps:} (1)~Choose and justify the edge stack (CDN, WAF mode,
TLS posture per traffic class) and the gateway/mesh split; draw the
request-path diagram for the payments path. (2)~Define the manifest
template (probes, resources policy, preStop, grace periods) with
per-service-class parameters. (3)~Define the progressive-delivery
standard: steps, gates ($\delta$, $\alpha$, $\beta$ per SLO tier),
rollback semantics, and the expand-and-contract pipeline stage that
blocks destructive migrations. (4)~Specify the long-lived-connection
strategy for WebSocket/gRPC (max connection age, rebalancing, drain
budgets). (5)~Write the incident runbook for ``canary gate refuses to
decide.'' (6)~Deliver as an architecture decision record set: one ADR
per load-bearing choice, each with the rejected alternatives.

\textbf{Acceptance Criteria:}
\begin{itemize}
  \item[$\square$] Every ADR names the rejected alternatives and the
        reversal cost.
  \item[$\square$] The payments path has an explicit TLS-termination
        decision with a compliance rationale (no plaintext last hop).
  \item[$\square$] The rollout standard includes per-tier gate
        parameters \emph{and} the evidence (simulator runs) that they
        detect the stated regression sizes.
  \item[$\square$] The WebSocket service's drain budget is derived from
        measured connection lifetimes, not defaults.
  \item[$\square$] A peer can run Capstone 17.3's linter against your
        template and find zero violations.
\end{itemize}

\textbf{Stretch Goals:} Add the multi-region failover half: health-based
DNS shifting, what ``rollback'' means when the failure is a region, and
where the canary gates live when traffic itself moves.

%% ----------------------------------------------------------------------------
\section{Python Dependencies Appendix}
\label{sec:ch17-deps}

All runnable code is offline and localhost-only after a one-time install.
From PowerShell:

\begin{minted}{text}
python -m venv .venv
.\.venv\Scripts\Activate.ps1
pip install -r requirements.txt
\end{minted}

\noindent\textbf{requirements.txt:}
\begin{minted}{text}
fastapi>=0.110
uvicorn>=0.29
httpx>=0.27
pytest>=8
pyyaml>=6.0
\end{minted}

\subsection{fastapi ($\geq$ 0.110)}
\textbf{Description:} The ASGI framework the Orders API is written in.
\textbf{Why included:} Listings 17.3--17.4 containerize and terminate a
FastAPI service; its middleware is the in-flight ledger the drain proof
relies on~\cite{fastapidocs}. \textbf{Alternatives:} Starlette directly
(the middleware pattern is raw ASGI and ports unchanged); Flask (WSGI ---
no async lifespan semantics, weaker fit for graceful-shutdown teaching).
\textbf{Installation:} \texttt{pip install "fastapi>=0.110"} in the
activated venv (PowerShell: keep the quotes so \texttt{>=} is not parsed
as redirection).

\subsection{uvicorn ($\geq$ 0.29)}
\textbf{Description:} The ASGI server whose signal handling is the
subject of Listing~\ref{lst:ch17-shutdown}: on \texttt{SIGTERM} (or
\texttt{SIGBREAK} via \texttt{CTRL\_BREAK\_EVENT} on Windows) it stops
accepting, waits for in-flight requests up to
\texttt{timeout\_graceful\_shutdown}, runs lifespan shutdown, and exits.
\textbf{Why included:} graceful shutdown cannot be taught as folklore ---
this chapter proves it against the real server. \textbf{Alternatives:}
hypercorn (comparable lifespan semantics); gunicorn with uvicorn workers
(the classic production shape whose own signal choreography --- master
vs.\ worker --- is an interview topic in itself).
\textbf{Installation:} \texttt{pip install "uvicorn>=0.29"}.

\subsection{httpx ($\geq$ 0.27)}
\textbf{Description:} Sync/async HTTP client. \textbf{Why included:} it
backs FastAPI's \texttt{TestClient} for any tests you write around the
listings, and is the natural client for extending the lab's load mix.
(The driver itself deliberately uses stdlib \texttt{urllib} so the proof
depends on nothing but Python.) \textbf{Alternatives:} \texttt{requests}
(synchronous only); raw \texttt{urllib} (what the driver uses --- proof
that you can, not that you should). \textbf{Installation:}
\texttt{pip install "httpx>=0.27"}.

\subsection{pytest ($\geq$ 8)}
\textbf{Description:} Test framework. \textbf{Why included:} the lab and
capstones wrap the demos as executable assertions (e.g., ``zero severed
requests''); exit codes from the demo driver map cleanly onto pytest
outcomes. \textbf{Alternatives:} \texttt{unittest} (stdlib; serviceable,
more ceremony). \textbf{Installation:} \texttt{pip install "pytest>=8"}.

\subsection{pyyaml ($\geq$ 6.0)}
\textbf{Description:} YAML parser. \textbf{Why included:} Capstone 17.3's
manifest linter parses Listing 17.2's YAML documents; parsing is the
honest way to review manifests --- grepping YAML for
\texttt{readOnlyRootFilesystem} misses indentation-level semantics that
\texttt{yaml.safe\_load} gets right. \textbf{Alternatives:}
\texttt{ruamel.yaml} (round-trips comments; heavier); \texttt{kubeconform}
/ \texttt{conftest} (real schema/policy validation when you graduate from
teaching lint rules). \textbf{Installation:}
\texttt{pip install "pyyaml>=6.0"}.

\subsection*{Toolchain (not pip --- install with winget)}
The chapter's reviewed-not-applied artifacts come alive with three tools,
all optional here: \textbf{Docker Desktop}
(\texttt{winget install Docker.DockerDesktop}) to build Listing 17.1 and
feel the layer cache; \textbf{kubectl}
(\texttt{winget install Kubernetes.kubectl}) plus a local cluster
(\texttt{kind} --- \texttt{winget install Kubernetes.kind} --- or Docker
Desktop's bundled Kubernetes) to actually apply Listing 17.2 and watch
probes and rollouts in \texttt{kubectl get pods -w}; and \textbf{k6}
(\texttt{winget install k6.k6}) to put real load on the canary endpoints
instead of synthetic fixtures. None is required: every behavior this
chapter claims is either proven by the runnable listings or stated as
platform semantics you will verify on first contact with a cluster.

%% ----------------------------------------------------------------------------
\section{Chapter Summary \& Key Takeaways}
\label{sec:ch17-summary}

\begin{itemize}
  \item A container image is a content-addressed layer stack: pin digests,
        split build from runtime, drop root, make the filesystem
        read-only, and treat image size as attack-surface accounting ---
        every line of Listing 17.1 is a security decision.
  \item Inside the container you are PID~1: exec-form entrypoints and
        real signal handlers are what make \texttt{SIGTERM} mean
        ``drain'' instead of ``die.'' The four-act sequence --- stop
        accepting, drain, timeout, exit --- is provable offline
        (Listings 17.3--17.4), and each act has a platform knob
        (readiness, graceful timeout, \texttt{terminationGracePeriodSeconds})
        that must agree with the others.
  \item The gateway is the policy enforcement point for north-south
        traffic: JWT offload, rate limiting, transformation, and canary
        routing configured once (Table~\ref{tab:ch17-gateway-matrix});
        ingress is its portable subset; the mesh owns east-west. Enforce
        each concern in exactly one tier.
  \item Probes are a recovery theory, not boilerplate: startup budgets
        boot time, readiness moves traffic without restarts, liveness
        restarts only process-local death --- and Scenario~B of the
        simulator shows the CrashLoopBackOff that results from conflating
        them.
  \item Requests schedule, limits throttle: CPU limits quantize latency
        into 100~ms stalls; memory limits OOM-kill. Set memory request
        $=$ limit, CPU request always, CPU limit deliberately.
  \item HPA scales on metrics against requests; stabilization windows and
        downscale policies are connection-protection, and PDBs are the
        arithmetic that keeps voluntary disruptions inside the SLO.
  \item Progressive delivery migrates traffic, not versions: size canary
        steps by requests, gate them with a sequential test with
        pre-registered $\delta$, $\alpha$, $\beta$
        (Listing~\ref{lst:ch17-canary}), and make rollback a pointer
        move --- which requires expand-and-contract schemas
        (Chapter~\ref{ch:versioning}).
  \item TLS placement is a visibility-versus-encryption decision per
        hop: terminate at the edge, re-encrypt inside, reserve
        passthrough for regulated payloads; automate every certificate;
        let the CDN cache only what is anonymous and keyed correctly.
  \item Termination is a concurrent race among SIGTERM, endpoint removal,
        and LB draining; the preStop sleep exists to rig it, and the
        missing preStop in Section~\ref{sec:ch17-case} is what 0.1\%
        dropped requests looks like from the inside.
\end{itemize}

%% ----------------------------------------------------------------------------
\section{Quick-Reference Card}
\label{sec:ch17-quickref}

\subsection*{Probe Behavior Table}
\begin{table}[h]
  \centering
  \small
  \begin{tabularx}{\textwidth}{l X X}
    \toprule
    \textbf{Probe} & \textbf{Failure consequence} & \textbf{Checks} \\
    \midrule
    Startup & container restarted after period $\times$ threshold budget & boot complete (cheap, no deps) \\
    Readiness & endpoint removed; re-added on first success; no restart & may touch critical deps, shallowly \\
    Liveness & container restarted in place after threshold streak & process-local only; never dependencies \\
    \bottomrule
  \end{tabularx}
\end{table}

Gating: startup runs alone until first success; readiness and liveness
begin only afterwards. Endpoint removal lags by seconds (kube-proxy, DNS,
LB convergence) --- rig the race with \texttt{preStop} sleep.

\subsection*{Rollout \& Autoscaling Cheat-Sheet}
\texttt{maxSurge: 1, maxUnavailable: 0} $\Rightarrow$ capacity never dips,
costs one surge pod $\bullet$ \texttt{maxUnavailable: 1, maxSurge: 0}
$\Rightarrow$ no extra quota, serves at $R{-}1$ windows $\bullet$
\texttt{progressDeadlineSeconds} fails wedged rollouts loudly $\bullet$
\texttt{revisionHistoryLimit: 5} keeps rollbacks instant $\bullet$
\texttt{terminationGracePeriodSeconds} $=$ preStop $+$ slowest request
$+$ margin $\bullet$ HPA: $\lceil \text{current} \times
\text{metric}/\text{target} \rceil$ against \emph{requests} $\bullet$
scaleDown window 300~s $+$ $\le$1 pod/60~s; scaleUp immediate $\bullet$
PDB: \texttt{minAvailable} $= R - 1$ for rollouts to proceed; never
${=}R$ (deadlocks drains); canary weights by pod ratio: canary pods $/$
total pods.

\subsection*{TLS Termination Matrix}
\begin{table}[h]
  \centering
  \small
  \begin{tabularx}{\textwidth}{l X X X}
    \toprule
    \textbf{Mode} & \textbf{Edge sees} & \textbf{Last hop} & \textbf{Use when} \\
    \midrule
    Terminate at edge & full L7 (WAF, JWT, routing) & plaintext & private network you truly control; legacy \\
    Re-encrypt & full L7 & TLS to backend (internal CA / mesh mTLS) & default for production APIs \\
    Passthrough & L4 only (no WAF, no L7 route) & end-to-end encrypted & regulated payloads; L7 policy unnecessary \\
    \bottomrule
  \end{tabularx}
\end{table}

Certificates: ACME (Let's Encrypt) at the edge; cert-manager or mesh
issuance inside; a manually renewed certificate is an outage with a
calendar date.

\section{Standardized Evidence Addendum}
\subsection{Benchmark Snapshot Template}
\begin{table}[h]
  \centering
  \small
  \begin{tabularx}{\textwidth}{l c c c c}
    \toprule
    \textbf{Config} & \textbf{p50 latency} & \textbf{p99 latency} & \textbf{Error rate} & \textbf{Throughput} \\
    \midrule
    Baseline & -- & -- & -- & 1.00x \\
    Candidate A & -- & -- & -- & -- \\
    Candidate B & -- & -- & -- & -- \\
    \bottomrule
  \end{tabularx}
  \caption{Minimum benchmark table to keep per chapter.}
\end{table}

\subsection{Request Trace Template}
\begin{itemize}
  \item \textbf{Request:} one representative API call for this chapter's topic.
  \item \textbf{Before:} behavior (headers, payload, latency, failure mode) with the naive approach.
  \item \textbf{After:} behavior with the technique in this chapter.
  \item \textbf{Result:} what changed in correctness, resilience, latency, and operability.
\end{itemize}

\subsection{Cost and Latency Budget Snapshot}
\begin{table}[h]
  \centering
  \small
  \begin{tabularx}{\textwidth}{l c c}
    \toprule
    \textbf{Stage} & \textbf{Latency budget} & \textbf{Cost driver} \\
    \midrule
    TLS / connection setup & -- & handshake round trips, keep-alive hit rate \\
    Request validation & -- & schema complexity, CPU per request \\
    Business logic and I/O & -- & downstream fan-out, query cost \\
    Serialization and egress & -- & payload size, compression ratio \\
    \bottomrule
  \end{tabularx}
  \caption{Operational budget template for this chapter's technique.}
\end{table}

\section{Configuration Preset Template}
\begin{minted}{text}
# api_policy.yaml
policy_version: "v1.0.0"
component: "<chapter_topic>"
timeout_ms: 0
retry_budget: 0
fallback_strategy: "fail_fast"
rollback_trigger: "error_budget_burn_gt_threshold"
\end{minted}

\section{CI Quality Gate Specification}
\begin{itemize}
  \item contract tests green against the pinned OpenAPI/AsyncAPI/Protobuf spec,
  \item no breaking schema changes without a version bump,
  \item error responses conform to RFC 9457 problem-details shape,
  \item benchmark deltas within approved regression bounds (latency and throughput),
  \item policy version and rollback plan present in release metadata.
\end{itemize}

\section{Incident Runbook Template}
\begin{enumerate}
  \item Detect regression from SLO burn-rate alerts or consumer reports.
  \item Scope impacted endpoints, versions, and consumer segments.
  \item Contain impact (rollback, feature flag, rate-limit tightening, or traffic shift).
  \item Diagnose with traces, structured logs, and contract diffs.
  \item Recover with validated fix and verified rollout procedure.
  \item Prevent recurrence via new regression tests and CI gates.
\end{enumerate}

\section{Cross-Chapter Decision Card}
\begin{table}[h]
  \centering
  \small
  \begin{tabularx}{\textwidth}{l X X}
    \toprule
    \textbf{Dimension} & \textbf{Choose this approach when} & \textbf{Prefer an alternative when} \\
    \midrule
    Use case & -- & -- \\
    Latency budget & -- & -- \\
    Operational cost & -- & -- \\
    Team maturity & -- & -- \\
    \bottomrule
  \end{tabularx}
  \caption{Decision card to complete per chapter.}
\end{table}

\section{ROI Model and Build-vs-Buy}
\begin{itemize}
  \item \textbf{Build when:} the capability is differentiating, compliance forbids vendors, or scale makes vendor pricing irrational.
  \item \textbf{Buy when:} the capability is commodity (auth, gateways, observability), time-to-market dominates, or staffing is thin.
  \item \textbf{Hidden costs to model:} on-call load, upgrade cadence, expertise ramp, data egress fees.
\end{itemize}

\section{Disaster Recovery Notes (RPO/RTO)}
\begin{itemize}
  \item Define recovery point objective (RPO): maximum acceptable data loss window.
  \item Define recovery time objective (RTO): maximum acceptable restoration time.
  \item Test restores regularly; an untested backup is not a backup.
\end{itemize}

